\documentclass[11pt,letterpaper]{article}

% ============================================================
% PACKAGES
% ============================================================

% Encoding and fonts
\usepackage[T1]{fontenc}
\usepackage[utf8]{inputenc}
\usepackage{lmodern}

% Page layout
\usepackage[margin=1in]{geometry}
\usepackage{setspace}
\usepackage{parskip}

% Mathematics
\usepackage{amsmath}
\usepackage{amssymb}
\usepackage{amsfonts}

% Tables and lists
\usepackage{booktabs}
\usepackage{tabularx}
\usepackage{longtable}
\usepackage{enumitem}

% Figures and graphics
\usepackage{graphicx}
\usepackage{xcolor}
\usepackage{float}

% Section formatting
\usepackage{titlesec}

% Code
\usepackage{listings}
\usepackage{upquote}

% Hyperlinks
\usepackage[
    colorlinks=true,
    linkcolor=blue!60!black,
    urlcolor=blue!60!black,
    citecolor=blue!60!black
]{hyperref}

% Header and footer
\usepackage{fancyhdr}
\usepackage{lastpage}


% ============================================================
% COLORS
% ============================================================

\definecolor{codegray}{gray}{0.45}
\definecolor{codebg}{gray}{0.96}


% ============================================================
% CODE LISTINGS
% ============================================================

\lstset{
    basicstyle=\ttfamily\small,
    backgroundcolor=\color{codebg},
    commentstyle=\color{codegray},
    numbers=none,
    frame=single,
    rulecolor=\color{black!20},
    breaklines=true,
    breakatwhitespace=false,
    showstringspaces=false,
    showspaces=false,
    showtabs=false,
    keepspaces=true,
    columns=fullflexible,
    upquote=true,
    tabsize=4,
    xleftmargin=0.5em,
    xrightmargin=0.5em,
    framexleftmargin=0.4em,
    framexrightmargin=0.4em
}


% ============================================================
% DOCUMENT FORMATTING
% ============================================================

\setstretch{1.15}

\titleformat{\section}
    {\Large\bfseries}
    {\thesection}{1em}{}

\titleformat{\subsection}
    {\large\bfseries}
    {\thesubsection}{1em}{}

\titleformat{\subsubsection}
    {\normalsize\bfseries}
    {\thesubsubsection}{1em}{}

% Show tools (subsections) in table of contents
\setcounter{tocdepth}{2}

% Number sections and subsections
\setcounter{secnumdepth}{2}


% ============================================================
% HEADER AND FOOTER
% ============================================================

\pagestyle{fancy}
\fancyhf{}

\fancyhead[L]{\small Tech Tools Cheatsheet}
\fancyhead[R]{\small Jonathan Ma}

\fancyfoot[C]{\thepage\ of \pageref{LastPage}}

\renewcommand{\headrulewidth}{0.4pt}


% ============================================================
% DOCUMENT
% ============================================================

\begin{document}


% ============================================================
% TITLE PAGE
% ============================================================

\begin{titlepage}
    \centering

    \vspace*{2.5cm}

    {\Huge\bfseries Tech Tools Cheatsheet\par}

    \vspace{0.5cm}

    {\Large A Reference Guide to Technical Tools\par}

    \vspace{1.5cm}

    {\large
    Created by
    \href{https://jonathanma03.github.io}{Jonathan Ma}
    \par}

    \vspace{0.5cm}

    {\large \today\par}

    \vfill

\end{titlepage}


% ============================================================
% TABLE OF CONTENTS
% ============================================================

\tableofcontents

\newpage


% ============================================================
% 1. DATA ACQUISITION
% ============================================================

\section{Data Acquisition}

% ============================================================
% REDCAP
% ============================================================

\subsection{REDCap}

\textbf{REDCap (Research Electronic Data Capture)} is a secure platform for
collecting and managing research data, particularly clinical and academic
research data.

The basic hierarchy is:

\[
\text{Project}
\rightarrow
\text{Record}
\rightarrow
\text{Instrument}
\rightarrow
\text{Field}
\]

For longitudinal studies:

\[
\text{Project}
\rightarrow
\text{Record}
\rightarrow
\text{Event}
\rightarrow
\text{Instrument}
\rightarrow
\text{Field}
\]

\textbf{Core terminology}

\begin{center}
\begin{tabularx}{\textwidth}{lX}
\toprule
Term & Meaning \\
\midrule
Project & A study, survey, registry, or research database \\
Record & One observational unit, usually a participant \\
Record ID & Unique identifier for a record \\
Instrument & A form containing related fields \\
Field & Individual variable collected in an instrument \\
Event & Study time point in a longitudinal project \\
Arm & Study pathway or group containing events \\
DAG & Data Access Group restricting records by site/group \\
Data Dictionary & CSV representation of fields and project structure \\
\bottomrule
\end{tabularx}
\end{center}

\textbf{Common field types}

\begin{itemize}[leftmargin=*]
    \item Text
    \item Notes
    \item Multiple choice / radio buttons
    \item Dropdown
    \item Checkbox
    \item Yes/No
    \item True/False
    \item Calculated field
    \item File upload
\end{itemize}

Text fields can also use validation such as integer, number, date, datetime,
email, and other formats.

\textbf{Variable names versus labels}

A field has both a machine-friendly variable name and a human-readable label.

For example:

\begin{lstlisting}
Variable: systolic_bp
Label:    Systolic Blood Pressure (mmHg)
\end{lstlisting}

Use stable, descriptive variable names. Avoid embedding information that may
later change into variable names.

Prefer:

\begin{lstlisting}
baseline_weight
treatment_group
visit_date
\end{lstlisting}

over ambiguous names such as:

\begin{lstlisting}
var1
question2
data3
\end{lstlisting}

\textbf{Branching logic}

Branching logic conditionally displays fields based on other responses.

\begin{lstlisting}
[smoker] = "1"
\end{lstlisting}

Multiple conditions:

\begin{lstlisting}
[smoker] = "1" and [age] >= 18
\end{lstlisting}

Alternative conditions:

\begin{lstlisting}
[treatment] = "1" or [treatment] = "2"
\end{lstlisting}

For checkbox fields, individual choices are referenced separately:

\begin{lstlisting}
[symptoms(1)] = "1"
\end{lstlisting}

\textbf{Calculated fields}

Calculated fields can derive values from existing fields.

For example, BMI:

\begin{lstlisting}
[weight_kg] / ([height_m] ^ 2)
\end{lstlisting}

Calculated fields are useful for simple deterministic quantities, but complex
analysis is generally better performed outside the data-entry layer.

\textbf{Longitudinal projects}

Longitudinal projects associate instruments with defined events such as:

\begin{lstlisting}
Baseline
Week 4
Week 8
Month 6
\end{lstlisting}

The same variable can therefore be observed at different study time points.

An event-specific reference can look like:

\begin{lstlisting}
[baseline_arm_1][weight_kg]
\end{lstlisting}

\textbf{Repeating instruments}

Repeating instruments allow multiple instances of the same form for a record,
such as repeated:

\begin{itemize}[leftmargin=*]
    \item adverse events;
    \item medications;
    \item hospital visits;
    \item follow-up assessments.
\end{itemize}

Do not confuse \textbf{repeating instruments} with \textbf{longitudinal
events}. Events represent planned study time points; repeating instruments
represent an arbitrary number of repeated instances.

\textbf{Data Dictionary}

The Data Dictionary is one of the most useful REDCap tools for larger
projects. It allows project fields to be inspected or defined in spreadsheet
form.

Important information includes:

\begin{itemize}[leftmargin=*]
    \item variable name;
    \item instrument/form name;
    \item field type;
    \item field label;
    \item choices/codings;
    \item validation;
    \item calculation;
    \item branching logic;
    \item required-field status.
\end{itemize}

For large projects, editing a Data Dictionary can be substantially faster and
more reproducible than manually creating every field through the interface.

\textbf{Coded values}

Always distinguish between the \textbf{stored value} and the
\textbf{displayed label}.

For example:

\begin{lstlisting}
0, Control
1, Treatment
\end{lstlisting}

The database may store:

\begin{lstlisting}
1
\end{lstlisting}

while the user sees:

\begin{lstlisting}
Treatment
\end{lstlisting}

This distinction matters when exporting and analyzing REDCap data.

\textbf{Checkboxes}

Checkboxes deserve particular attention because multiple choices can be
selected simultaneously. Exports may represent each option as a separate
indicator variable rather than one categorical variable.

Conceptually:

\begin{lstlisting}
symptoms___1
symptoms___2
symptoms___3
\end{lstlisting}

These may correspond to separate yes/no indicators for each response option.

\textbf{Data quality}

Useful REDCap design practices include:

\begin{itemize}[leftmargin=*]
    \item validate values at entry when possible;
    \item use required fields selectively;
    \item define sensible minimum/maximum ranges;
    \item use branching logic to prevent irrelevant questions;
    \item avoid unnecessary free-text fields;
    \item use standardized coding consistently;
    \item test instruments before production use.
\end{itemize}

\textbf{Practical tips}

\begin{itemize}[leftmargin=*]
    \item Build a small prototype before constructing a large project.
    \item Decide variable naming conventions before adding hundreds of fields.
    \item Keep a copy of the Data Dictionary for reproducibility.
    \item Check whether exports use raw coded values or displayed labels.
    \item Be especially careful with checkbox exports and repeating instruments.
    \item Test branching logic using several realistic records.
    \item Avoid changing variable meaning after data collection has started.
    \item When possible, constrain data during collection instead of cleaning
    preventable errors afterward.
\end{itemize}

\textbf{Interview notes}

Be able to explain the distinction between a record, instrument, field, event,
and repeating instrument. A useful design principle is that REDCap should
enforce as much valid structure as practical \emph{during collection}, because
preventing invalid data is generally preferable to repairing it downstream.


% ============================================================
% FHIR
% ============================================================

\subsection{FHIR}

\textbf{FHIR (Fast Healthcare Interoperability Resources)} is an HL7 standard
for representing and exchanging healthcare information.

FHIR organizes healthcare information into modular
\textbf{resources}.

\textbf{Common resources}

\begin{center}
\begin{tabularx}{\textwidth}{lX}
\toprule
Resource & Represents \\
\midrule
Patient & Demographics and patient identity \\
Encounter & Interaction with a healthcare system \\
Condition & Diagnosis or clinical condition \\
Observation & Measurement, lab result, or vital sign \\
Procedure & Clinical procedure \\
Medication & Medication definition \\
MedicationRequest & Medication order or prescription \\
DiagnosticReport & Diagnostic/laboratory report \\
Practitioner & Healthcare professional \\
Organization & Healthcare organization \\
AllergyIntolerance & Allergy or intolerance \\
Immunization & Vaccination event \\
\bottomrule
\end{tabularx}
\end{center}

A useful mental model is:

\[
\text{Patient}
\longleftarrow
\begin{cases}
\text{Encounter} \\
\text{Observation} \\
\text{Condition} \\
\text{Procedure} \\
\text{MedicationRequest}
\end{cases}
\]

FHIR therefore does \emph{not} generally represent an entire patient's
history as one giant record.

\textbf{Resource structure}

FHIR resources are commonly represented using JSON or XML.

A simplified Patient resource:

\begin{lstlisting}
{
  "resourceType": "Patient",
  "id": "123",
  "name": [
    {
      "family": "Smith",
      "given": ["Jane"]
    }
  ],
  "gender": "female",
  "birthDate": "1985-04-12"
}
\end{lstlisting}

The \texttt{resourceType} identifies what kind of resource the object
represents.

\textbf{References}

FHIR resources are connected through references.

An Observation might contain:

\begin{lstlisting}
"subject": {
  "reference": "Patient/123"
}
\end{lstlisting}

An Encounter might similarly reference the patient:

\begin{lstlisting}
"subject": {
  "reference": "Patient/123"
}
\end{lstlisting}

This means that interpreting clinical data often requires understanding
relationships among several resources.

\textbf{Observation}

Observation is one of the most important resources for data work.

A simplified blood-pressure observation might contain:

\begin{lstlisting}
{
  "resourceType": "Observation",
  "status": "final",
  "code": {
    "coding": [
      {
        "system": "http://loinc.org",
        "code": "85354-9"
      }
    ]
  },
  "subject": {
    "reference": "Patient/123"
  }
}
\end{lstlisting}

An Observation is not synonymous with a laboratory result. It can represent
many measured clinical facts, including:

\begin{itemize}[leftmargin=*]
    \item laboratory measurements;
    \item vital signs;
    \item body measurements;
    \item clinical scores;
    \item other measured observations.
\end{itemize}

\textbf{Clinical coding systems}

FHIR resources frequently contain codes from external clinical terminology
systems.

\begin{center}
\begin{tabularx}{\textwidth}{lX}
\toprule
System & Typical Use \\
\midrule
LOINC & Laboratory tests and clinical observations \\
SNOMED CT & Detailed clinical terminology \\
ICD-10 & Diagnosis classification \\
RxNorm & Medication terminology \\
CPT & Procedures/services, particularly U.S. billing contexts \\
\bottomrule
\end{tabularx}
\end{center}

A coded element often contains:

\begin{lstlisting}
{
  "system": "http://loinc.org",
  "code": "8480-6",
  "display": "Systolic blood pressure"
}
\end{lstlisting}

Keep the distinction clear:

\[
\text{FHIR}
=
\text{structure/interoperability}
\]

while systems such as LOINC and SNOMED CT provide much of the
\text{clinical terminology}.

\textbf{CodeableConcept}

FHIR frequently represents coded clinical information using a
\texttt{CodeableConcept}.

A concept can contain multiple codings:

\begin{lstlisting}
"code": {
  "coding": [
    {
      "system": "...",
      "code": "...",
      "display": "..."
    }
  ],
  "text": "Human-readable description"
}
\end{lstlisting}

Do not assume the human-readable \texttt{display} string is the canonical
identifier. The coding system and code carry the standardized meaning.

\textbf{Cardinality}

FHIR specifications use cardinalities such as:

\begin{lstlisting}
0..1
1..1
0..*
1..*
\end{lstlisting}

Interpretation:

\begin{center}
\begin{tabular}{ll}
\toprule
Cardinality & Meaning \\
\midrule
\texttt{0..1} & Optional, at most one \\
\texttt{1..1} & Required, exactly one \\
\texttt{0..*} & Optional, any number \\
\texttt{1..*} & At least one \\
\bottomrule
\end{tabular}
\end{center}

This is useful when reading FHIR specifications and understanding why some
fields are arrays while others are singular values.

\textbf{Profiles}

FHIR itself is flexible. A \textbf{profile} constrains a base resource for a
particular use case.

A profile may:

\begin{itemize}[leftmargin=*]
    \item require fields that are optional in base FHIR;
    \item restrict allowable terminology;
    \item constrain cardinality;
    \item define expected extensions.
\end{itemize}

Therefore, ``FHIR compliant'' does not imply that every implementation stores
identical fields in identical ways.

\textbf{Extensions}

FHIR supports extensions for information not represented directly by the base
resource.

Conceptually:

\begin{lstlisting}
"extension": [
  {
    "url": "...",
    "valueString": "..."
  }
]
\end{lstlisting}

Extensions are an intentional part of FHIR's design, not necessarily a
violation of the standard.

\textbf{Bundles}

Collections of resources can be represented as a \textbf{Bundle}.

\begin{lstlisting}
{
  "resourceType": "Bundle",
  "type": "collection",
  "entry": [
    {
      "resource": { ... }
    },
    {
      "resource": { ... }
    }
  ]
}
\end{lstlisting}

The important practical point is that a Bundle contains multiple resource
entries; it is not itself a Patient, Observation, or Encounter.

\textbf{Practical tips}

\begin{itemize}[leftmargin=*]
    \item Always inspect \texttt{resourceType} before assuming an object's structure.
    \item Expect nested objects and arrays rather than flat tabular data.
    \item Do not identify clinical concepts using display text alone.
    \item Preserve coding-system information along with clinical codes.
    \item Expect resources to reference other resources.
    \item Check the FHIR version and implementation profile being used.
    \item Do not assume an optional field will be populated simply because the
    standard defines it.
    \item Be careful when flattening FHIR into tables; one-to-many relationships can
    easily duplicate patient-level rows.
\end{itemize}

\textbf{Interview notes}

The most important distinction is:

\[
\boxed{\text{FHIR is an interoperability standard, not an EHR or analytics schema.}}
\]

Be comfortable explaining resources, references, terminology systems,
profiles, extensions, and cardinality. Also understand why converting nested
FHIR resources into analytical tables is not always a simple one-row-per-patient
operation.


% ============================================================
% EPIC
% ============================================================

\subsection{Epic}

\textbf{Epic} is a major Electronic Health Record (EHR) platform used for
clinical documentation, orders, results, scheduling, billing, and other
healthcare workflows.

For technical work, the most important point is that ``Epic'' refers to an
ecosystem rather than a single flat database.

\textbf{Useful conceptual layers}

\begin{center}
\begin{tabularx}{\textwidth}{lX}
\toprule
Component & Primary Role \\
\midrule
Epic EHR & Operational clinical system \\
Hyperspace & Traditional Epic end-user application \\
Hyperdrive & Web-based Epic client environment \\
Reporting Workbench & Operational reporting \\
Clarity & Relational reporting database \\
Caboodle & Enterprise analytics/data warehouse \\
\bottomrule
\end{tabularx}
\end{center}

A useful mental model is:

\[
\text{Clinical Workflow}
\rightarrow
\text{Operational Epic Data}
\rightarrow
\begin{cases}
\text{Operational Reporting} \\
\text{Clarity} \\
\text{Caboodle}
\end{cases}
\]

\textbf{Operational versus analytical data}

The live EHR is optimized for clinical operations:

\begin{itemize}[leftmargin=*]
    \item recording encounters;
    \item documenting clinical notes;
    \item entering orders;
    \item recording results;
    \item scheduling;
    \item billing and administrative workflows.
\end{itemize}

Large analytical workloads generally rely on downstream reporting or
warehouse systems rather than treating the live operational system as a
general-purpose analytics database.

\textbf{Clarity}

\textbf{Clarity} is Epic's relational reporting database. Operational Epic
data is transformed into relational tables suitable for reporting, research,
and analysis.

Conceptually, analysts may encounter data corresponding to:

\begin{itemize}[leftmargin=*]
    \item patients;
    \item encounters;
    \item diagnoses;
    \item medications;
    \item orders;
    \item procedures;
    \item laboratory results;
    \item providers.
\end{itemize}

Clarity is important because EHR data is naturally relational. A patient can
have many encounters, and each encounter can have many diagnoses, orders,
medications, or results.

\[
\text{Patient}
\rightarrow
\text{Encounter}
\rightarrow
\begin{cases}
\text{Diagnosis} \\
\text{Order} \\
\text{Medication} \\
\text{Result}
\end{cases}
\]

\textbf{Caboodle}

\textbf{Caboodle} is Epic's enterprise data warehouse. It provides a more
analytics-oriented representation of data from Epic and potentially other
organizational systems.

A useful distinction is:

\[
\text{Clarity}
\approx
\text{detailed relational reporting data}
\]

\[
\text{Caboodle}
\approx
\text{enterprise analytics warehouse}
\]

Exact implementations and available data vary across organizations.

\textbf{Reporting Workbench}

Reporting Workbench supports operational reporting closer to day-to-day EHR
workflows. It is useful for operational lists and reports but serves a
different purpose from a large analytical warehouse.

\textbf{Encounter-centric data}

One of the most important concepts when working with EHR data is that
\textbf{patient} and \textbf{encounter} are different units.

A patient may have:

\[
1 \text{ patient}
\rightarrow
N \text{ encounters}
\]

and one encounter may have:

\[
1 \text{ encounter}
\rightarrow
N \text{ diagnoses/results/orders}
\]

Joining these tables without understanding their grain can dramatically
multiply rows.

Before using an EHR dataset, ask:

\begin{quote}
What does one row represent?
\end{quote}

Possible answers include:

\begin{itemize}[leftmargin=*]
    \item one patient;
    \item one encounter;
    \item one diagnosis;
    \item one medication order;
    \item one lab result;
    \item one procedure.
\end{itemize}

\textbf{Dates and timestamps}

Clinical data can contain several timestamps describing different events.

For example, a laboratory record might involve:

\[
\text{Order Time}
\rightarrow
\text{Collection Time}
\rightarrow
\text{Result Time}
\]

These are not interchangeable.

Similarly, an encounter may have scheduled, arrival, admission, discharge,
and documentation timestamps.

Always determine which timestamp corresponds to the scientific or operational
question.

\textbf{Clinical codes}

EHR data can contain both standardized and organization-specific codes.

Do not assume that:

\begin{itemize}[leftmargin=*]
    \item every diagnosis is represented identically across institutions;
    \item every local laboratory code maps trivially to one standard code;
    \item displayed names are stable identifiers;
    \item absence of a code necessarily means absence of the condition.
\end{itemize}

\textbf{Practical tips}

\begin{itemize}[leftmargin=*]
    \item Determine the grain of every table before joining.
    \item Check key uniqueness before assuming a one-to-one relationship.
    \item Distinguish patient-level identifiers from encounter-level identifiers.
    \item Inspect date semantics rather than selecting whichever timestamp is easiest.
    \item Expect repeated measurements and repeated encounters.
    \item Treat missing clinical data carefully; ``not recorded'' is not always
    equivalent to ``did not occur.''
    \item Preserve provenance when combining data from multiple clinical sources.
    \item Remember that Epic implementations differ across healthcare organizations.
\end{itemize}

\textbf{Interview notes}

Be able to explain why operational EHR systems and analytical warehouses have
different purposes. For data roles, emphasize \textbf{grain, temporal
semantics, join cardinality, provenance, and clinical coding}. A technically
correct join can still produce a clinically incorrect dataset if these are not
understood.


% ============================================================
% BLOOMBERG TERMINAL
% ============================================================

\subsection{Bloomberg Terminal}

\textbf{Bloomberg Terminal} is a professional financial information platform
used for market data, company fundamentals, economic data, news, screening,
valuation, and financial analysis.

Bloomberg workflows rely heavily on short mnemonic commands followed by the
\texttt{<GO>} key.

\textbf{Security syntax}

A security is commonly specified using:

\[
\text{Identifier}
+
\text{Market/Asset Type}
\]

Examples:

\begin{lstlisting}
AAPL US Equity <GO>
NVDA US Equity <GO>
SPX Index <GO>
EURUSD Curncy <GO>
\end{lstlisting}

Then apply a function:

\begin{lstlisting}
AAPL US Equity DES <GO>
AAPL US Equity GP <GO>
AAPL US Equity FA <GO>
\end{lstlisting}

\textbf{Useful everyday functions}

\begin{center}
\begin{tabularx}{\textwidth}{lX}
\toprule
Function & Purpose \\
\midrule
DES & Security description / overview \\
GP & Price chart \\
HP & Historical prices \\
FA & Financial analysis \\
DVD & Dividend information \\
CN & Company news \\
ANR & Analyst recommendations \\
ERN & Earnings information \\
RV & Relative valuation \\
EQS & Equity screening \\
WEI & Global equity indices \\
ECST & Economic statistics \\
TOP & Top Bloomberg news \\
HELP & Context-sensitive help \\
\bottomrule
\end{tabularx}
\end{center}

A good practical habit is to use:

\begin{lstlisting}
DES <GO>
\end{lstlisting}

when unfamiliar with a security, then move into more specialized functions.

\textbf{Asset-class identifiers}

Common Bloomberg security types include:

\begin{center}
\begin{tabular}{ll}
\toprule
Suffix & Asset Type \\
\midrule
Equity & Equity security \\
Index & Index/rate series \\
Curncy & Currency \\
Corp & Corporate bond \\
Govt & Government security \\
Comdty & Commodity \\
\bottomrule
\end{tabular}
\end{center}

Do not assume a ticker alone uniquely identifies a financial instrument.

\textbf{Excel Add-In}

The Bloomberg Excel Add-In is extremely useful for bringing Terminal data
into spreadsheet workflows.

Three functions are especially important:

\[
\boxed{\texttt{BDP},\quad \texttt{BDH},\quad \texttt{BDS}}
\]

\textbf{BDP: Bloomberg Data Point}

Retrieve a current/reference value:

\begin{lstlisting}
=BDP("AAPL US Equity","PX_LAST")
\end{lstlisting}

Examples:

\begin{lstlisting}
=BDP("AAPL US Equity","CUR_MKT_CAP")
=BDP("AAPL US Equity","PE_RATIO")
\end{lstlisting}

\textbf{BDH: Bloomberg Data History}

Retrieve historical time-series data:

\begin{lstlisting}
=BDH(
    "AAPL US Equity",
    "PX_LAST",
    "01/01/2026",
    "10/01/2026"
)
\end{lstlisting}

\textbf{BDS: Bloomberg Data Set}

Retrieve bulk/multi-row information:

\begin{lstlisting}
=BDS(
    "AAPL US Equity",
    "TOP_20_HOLDERS_PUBLIC_FILINGS"
)
\end{lstlisting}

Remember:

\begin{center}
\begin{tabular}{ll}
\toprule
Function & Use \\
\midrule
BDP & One/reference data point \\
BDH & Historical time series \\
BDS & Bulk/multi-row data \\
\bottomrule
\end{tabular}
\end{center}

\textbf{Common fields}

Bloomberg data fields use mnemonics such as:

\begin{lstlisting}
PX_LAST
PX_OPEN
PX_HIGH
PX_LOW
VOLUME
CUR_MKT_CAP
PE_RATIO
\end{lstlisting}

When unsure of a field, use Bloomberg's field search rather than guessing the
mnemonic.

\textbf{Historical data}

Before using downloaded financial time series, verify:

\begin{itemize}[leftmargin=*]
    \item observation frequency;
    \item date range;
    \item currency;
    \item trading calendar;
    \item missing/non-trading dates;
    \item corporate-action treatment;
    \item adjusted versus unadjusted values.
\end{itemize}

\textbf{Corporate actions}

Stock splits, dividends, spin-offs, and other corporate actions can affect
historical prices and returns.

For example, a raw price discontinuity caused by a stock split is not an
economic return:

\[
P_t \ll P_{t-1}
\]

may reflect a mechanical adjustment rather than a market loss.

Always know whether the requested series is adjusted appropriately for the
analysis.

\textbf{Point-in-time data}

For historical modeling and backtesting, distinguish:

\[
\text{Current Database Value}
\neq
\text{Value Known at Historical Date}
\]

Financial databases may contain revised fundamentals, classifications, or
other information that was not available to market participants at the time.

Using current information in a historical model creates
\textbf{look-ahead bias}.

Examples of potential problems include:

\begin{itemize}[leftmargin=*]
    \item revised economic data;
    \item restated financial statements;
    \item current index constituents;
    \item current industry classifications;
    \item later analyst estimates.
\end{itemize}

\textbf{Survivorship bias}

A current list of securities excludes firms that disappeared through
bankruptcy, acquisition, delisting, or other events.

Backtesting only currently surviving firms can therefore produce overly
optimistic results.

\textbf{Frequency alignment}

When combining Bloomberg series, frequencies may differ:

\[
\text{Daily Equity Prices}
+
\text{Monthly CPI}
+
\text{Quarterly Earnings}
\]

Do not mechanically forward-fill or merge series without considering when
information actually became available.

For macroeconomic data in particular:

\[
\text{Reference Period}
\neq
\text{Release Date}
\]

A value describing September may not have been known until October.

\textbf{Practical tips}

\begin{itemize}[leftmargin=*]
    \item Start with \texttt{DES} when exploring an unfamiliar security.
    \item Remember \texttt{BDP = point}, \texttt{BDH = history},
    \texttt{BDS = dataset}.
    \item Confirm security type when tickers are ambiguous.
    \item Verify units and currencies before comparing securities.
    \item Check whether historical prices account for corporate actions.
    \item Record the field mnemonic used to make an extraction reproducible.
    \item For backtests, ask whether the data are truly point-in-time.
    \item Check index membership historically rather than using only current
    constituents.
    \item Treat release dates carefully when combining market and macroeconomic data.
\end{itemize}

\textbf{Interview notes}

For finance/data interviews, the most important practical issues are often not
how to click through Bloomberg, but whether you understand the data being
retrieved. Be ready to discuss \textbf{adjusted prices, corporate actions,
point-in-time data, look-ahead bias, survivorship bias, security identifiers,
and frequency alignment}.

For Bloomberg-specific usage, know the distinction between
\texttt{BDP}, \texttt{BDH}, and \texttt{BDS}, and recognize common Terminal
functions such as \texttt{DES}, \texttt{GP}, \texttt{FA}, \texttt{EQS}, and
\texttt{ECST}.


% ============================================================
% 2. STORAGE AND QUERY
% ============================================================

\section{Storage and Query}

% ============================================================
% SQL
% ============================================================

\subsection{SQL}

\textbf{SQL (Structured Query Language)} is the standard language for working
with relational databases. For day-to-day analytics, the most important skills
are selecting and filtering data, joining tables, aggregating observations,
using window functions, and understanding the grain of a query result.

A useful mental model for a SQL query is:

\[
\texttt{FROM}
\rightarrow
\texttt{WHERE}
\rightarrow
\texttt{GROUP BY}
\rightarrow
\texttt{HAVING}
\rightarrow
\texttt{SELECT}
\rightarrow
\texttt{ORDER BY}
\rightarrow
\texttt{LIMIT}
\]

This is the approximate \emph{logical} processing order, not necessarily the
physical execution order chosen by the database optimizer.


\subsubsection{Core SQL}

\textbf{Basic query}

\begin{lstlisting}[language=SQL]
SELECT
    customer_id,
    revenue
FROM customers
WHERE revenue > 1000
ORDER BY revenue DESC
LIMIT 10;
\end{lstlisting}

Avoid using:

\begin{lstlisting}[language=SQL]
SELECT *
\end{lstlisting}

in production queries when only a few columns are required. Explicit columns
make the expected schema clearer and can reduce unnecessary data reads.

\textbf{Filtering}

\begin{lstlisting}[language=SQL]
SELECT *
FROM employees
WHERE salary >= 80000
  AND department = 'Analytics';
\end{lstlisting}

Common conditions:

\begin{lstlisting}[language=SQL]
WHERE x = 10
WHERE x <> 10
WHERE x > 10
WHERE x BETWEEN 10 AND 20
WHERE x IN (10, 20, 30)
WHERE name LIKE 'Jon%'
WHERE name LIKE '%Smith%'
WHERE x IS NULL
WHERE x IS NOT NULL
\end{lstlisting}

\textbf{NULL}

\texttt{NULL} represents an unknown or missing value.

Do not use:

\begin{lstlisting}[language=SQL]
WHERE x = NULL
\end{lstlisting}

Use:

\begin{lstlisting}[language=SQL]
WHERE x IS NULL
\end{lstlisting}

Similarly:

\begin{lstlisting}[language=SQL]
WHERE x IS NOT NULL
\end{lstlisting}

SQL uses three-valued logic:

\[
\text{TRUE},\qquad \text{FALSE},\qquad \text{UNKNOWN}
\]

Comparisons involving \texttt{NULL} generally evaluate to
\texttt{UNKNOWN}.

\textbf{COALESCE}

Return the first non-null value:

\begin{lstlisting}[language=SQL]
SELECT
    customer_id,
    COALESCE(phone, mobile_phone, 'Unknown') AS phone
FROM customers;
\end{lstlisting}

Replace missing numeric values:

\begin{lstlisting}[language=SQL]
SELECT
    COALESCE(revenue, 0) AS revenue
FROM customers;
\end{lstlisting}

Be careful: replacing missing values with zero changes their meaning and
should only be done when zero is scientifically or operationally appropriate.

\textbf{CASE}

\begin{lstlisting}[language=SQL]
SELECT
    customer_id,
    revenue,
    CASE
        WHEN revenue >= 10000 THEN 'High'
        WHEN revenue >= 5000 THEN 'Medium'
        ELSE 'Low'
    END AS revenue_group
FROM customers;
\end{lstlisting}

Conditional indicators are also common:

\begin{lstlisting}[language=SQL]
CASE
    WHEN age >= 18 THEN 1
    ELSE 0
END AS adult
\end{lstlisting}

\textbf{Aggregation}

Common aggregate functions:

\begin{lstlisting}[language=SQL]
COUNT(*)
COUNT(column_name)
COUNT(DISTINCT column_name)

SUM(x)
AVG(x)
MIN(x)
MAX(x)
\end{lstlisting}

Example:

\begin{lstlisting}[language=SQL]
SELECT
    department,
    COUNT(*) AS n,
    AVG(salary) AS avg_salary
FROM employees
GROUP BY department;
\end{lstlisting}

A crucial distinction:

\begin{lstlisting}[language=SQL]
COUNT(*)
\end{lstlisting}

counts rows, while:

\begin{lstlisting}[language=SQL]
COUNT(column_name)
\end{lstlisting}

counts non-null values in that column.

\textbf{WHERE versus HAVING}

\texttt{WHERE} filters rows before aggregation:

\begin{lstlisting}[language=SQL]
SELECT
    department,
    AVG(salary) AS avg_salary
FROM employees
WHERE active = 1
GROUP BY department;
\end{lstlisting}

\texttt{HAVING} filters groups after aggregation:

\begin{lstlisting}[language=SQL]
SELECT
    department,
    AVG(salary) AS avg_salary
FROM employees
GROUP BY department
HAVING AVG(salary) > 90000;
\end{lstlisting}

A useful rule is:

\[
\texttt{WHERE}
=
\text{filter observations}
\]

\[
\texttt{HAVING}
=
\text{filter aggregated groups}
\]


\textbf{Joins}

\begin{center}
\begin{tabularx}{\textwidth}{lX}
\toprule
Join & Result \\
\midrule
INNER JOIN & Only rows with matches in both tables \\
LEFT JOIN & Every left-table row plus matching right-table data \\
RIGHT JOIN & Every right-table row plus matching left-table data \\
FULL OUTER JOIN & All rows from both tables \\
CROSS JOIN & Cartesian product of both tables \\
\bottomrule
\end{tabularx}
\end{center}

Example:

\begin{lstlisting}[language=SQL]
SELECT
    c.customer_id,
    c.name,
    o.order_id,
    o.amount
FROM customers AS c
LEFT JOIN orders AS o
    ON c.customer_id = o.customer_id;
\end{lstlisting}

The most important practical question before joining is:

\begin{quote}
What is the relationship between these keys?
\end{quote}

Possible relationships are:

\[
1:1,\qquad 1:N,\qquad N:1,\qquad N:M
\]

If one customer has five orders, joining one customer row to the order table
creates five output rows.

Before trusting a join, check key uniqueness.

For example:

\begin{lstlisting}[language=SQL]
SELECT
    customer_id,
    COUNT(*) AS n
FROM customers
GROUP BY customer_id
HAVING COUNT(*) > 1;
\end{lstlisting}

If this returns rows, \texttt{customer\_id} is not unique in that table.

\textbf{Duplicate rows after joins}

Suppose:

\[
\text{Patient}
\rightarrow
5\text{ diagnoses}
\]

and:

\[
\text{Patient}
\rightarrow
10\text{ medications}
\]

Joining both directly by patient can potentially produce:

\[
5 \times 10 = 50
\]

rows for the patient.

This is a common real-world source of incorrect analyses.

\textbf{Subqueries}

\begin{lstlisting}[language=SQL]
SELECT *
FROM employees
WHERE salary > (
    SELECT AVG(salary)
    FROM employees
);
\end{lstlisting}

Subqueries can also appear in \texttt{FROM} or be used with operators such as
\texttt{EXISTS}.

\textbf{EXISTS}

\begin{lstlisting}[language=SQL]
SELECT *
FROM customers AS c
WHERE EXISTS (
    SELECT 1
    FROM orders AS o
    WHERE o.customer_id = c.customer_id
);
\end{lstlisting}

This asks whether at least one matching order exists without requiring the
matching rows to be returned.

\textbf{Common Table Expressions}

CTEs make multi-stage queries easier to read.

\begin{lstlisting}[language=SQL]
WITH customer_totals AS (
    SELECT
        customer_id,
        SUM(amount) AS total_spend
    FROM orders
    GROUP BY customer_id
)

SELECT *
FROM customer_totals
WHERE total_spend > 10000;
\end{lstlisting}

Multiple CTEs:

\begin{lstlisting}[language=SQL]
WITH recent_orders AS (
    SELECT *
    FROM orders
    WHERE order_date >= '2026-01-01'
),
customer_totals AS (
    SELECT
        customer_id,
        SUM(amount) AS total_spend
    FROM recent_orders
    GROUP BY customer_id
)

SELECT *
FROM customer_totals;
\end{lstlisting}

CTEs are especially useful for breaking complicated transformations into
logical stages.

\textbf{Window functions}

Window functions calculate values across related rows without collapsing the
result to one row per group.

General structure:

\begin{lstlisting}[language=SQL]
function(...) OVER (
    PARTITION BY ...
    ORDER BY ...
)
\end{lstlisting}

For example:

\begin{lstlisting}[language=SQL]
SELECT
    employee_id,
    department,
    salary,
    AVG(salary) OVER (
        PARTITION BY department
    ) AS department_avg
FROM employees;
\end{lstlisting}

Unlike \texttt{GROUP BY}, every employee remains in the output.

\textbf{ROW\_NUMBER, RANK, and DENSE\_RANK}

\begin{lstlisting}[language=SQL]
SELECT
    employee_id,
    department,
    salary,
    ROW_NUMBER() OVER (
        PARTITION BY department
        ORDER BY salary DESC
    ) AS rn
FROM employees;
\end{lstlisting}

Differences:

\begin{center}
\begin{tabularx}{\textwidth}{lX}
\toprule
Function & Behavior \\
\midrule
ROW\_NUMBER() & Unique sequential number for every row \\
RANK() & Ties share a rank and subsequent ranks contain gaps \\
DENSE\_RANK() & Ties share a rank without subsequent gaps \\
\bottomrule
\end{tabularx}
\end{center}

If salaries are:

\begin{lstlisting}
100
100
90
\end{lstlisting}

then:

\begin{center}
\begin{tabular}{rrrr}
\toprule
Salary & ROW\_NUMBER & RANK & DENSE\_RANK \\
\midrule
100 & 1 & 1 & 1 \\
100 & 2 & 1 & 1 \\
90  & 3 & 3 & 2 \\
\bottomrule
\end{tabular}
\end{center}

\textbf{Top $N$ per group}

A very common interview pattern:

\begin{lstlisting}[language=SQL]
WITH ranked AS (
    SELECT
        employee_id,
        department,
        salary,
        ROW_NUMBER() OVER (
            PARTITION BY department
            ORDER BY salary DESC
        ) AS rn
    FROM employees
)

SELECT *
FROM ranked
WHERE rn <= 3;
\end{lstlisting}

\textbf{LAG and LEAD}

Access previous or subsequent rows:

\begin{lstlisting}[language=SQL]
SELECT
    date,
    price,
    LAG(price) OVER (
        ORDER BY date
    ) AS previous_price
FROM prices;
\end{lstlisting}

Compute a change:

\begin{lstlisting}[language=SQL]
WITH lagged AS (
    SELECT
        date,
        price,
        LAG(price) OVER (
            ORDER BY date
        ) AS previous_price
    FROM prices
)

SELECT
    date,
    price,
    price - previous_price AS price_change
FROM lagged;
\end{lstlisting}

\textbf{Running totals}

\begin{lstlisting}[language=SQL]
SELECT
    date,
    revenue,
    SUM(revenue) OVER (
        ORDER BY date
        ROWS BETWEEN UNBOUNDED PRECEDING
                 AND CURRENT ROW
    ) AS cumulative_revenue
FROM daily_sales;
\end{lstlisting}

\textbf{Deduplication}

Find duplicate keys:

\begin{lstlisting}[language=SQL]
SELECT
    customer_id,
    COUNT(*) AS n
FROM customers
GROUP BY customer_id
HAVING COUNT(*) > 1;
\end{lstlisting}

Keep the most recent row per customer:

\begin{lstlisting}[language=SQL]
WITH ranked AS (
    SELECT
        *,
        ROW_NUMBER() OVER (
            PARTITION BY customer_id
            ORDER BY updated_at DESC
        ) AS rn
    FROM customers
)

SELECT *
FROM ranked
WHERE rn = 1;
\end{lstlisting}

\textbf{Set operations}

\begin{lstlisting}[language=SQL]
SELECT customer_id
FROM online_customers

UNION

SELECT customer_id
FROM store_customers;
\end{lstlisting}

Important operations:

\begin{itemize}[leftmargin=*]
    \item \texttt{UNION}: combines results and removes duplicates;
    \item \texttt{UNION ALL}: combines results and preserves duplicates;
    \item \texttt{INTERSECT}: rows appearing in both queries;
    \item \texttt{EXCEPT}: rows in the first query but not the second.
\end{itemize}

Use \texttt{UNION ALL} when duplicate removal is unnecessary.

\textbf{Data modification}

Insert:

\begin{lstlisting}[language=SQL]
INSERT INTO employees (
    employee_id,
    name,
    salary
)
VALUES (
    101,
    'Alice',
    95000
);
\end{lstlisting}

Update:

\begin{lstlisting}[language=SQL]
UPDATE employees
SET salary = 100000
WHERE employee_id = 101;
\end{lstlisting}

Delete:

\begin{lstlisting}[language=SQL]
DELETE FROM employees
WHERE employee_id = 101;
\end{lstlisting}

Always double-check the \texttt{WHERE} clause before running
\texttt{UPDATE} or \texttt{DELETE}.

A useful habit is to first run:

\begin{lstlisting}[language=SQL]
SELECT *
FROM employees
WHERE employee_id = 101;
\end{lstlisting}

and verify the affected rows before modifying them.

\textbf{Keys}

A \textbf{primary key} uniquely identifies a row.

A \textbf{foreign key} references a key in another table.

For example:

\[
\texttt{customers.customer\_id}
\longleftarrow
\texttt{orders.customer\_id}
\]

Primary and foreign keys encode relationships between tables.

\textbf{Indexes}

Indexes can accelerate lookups, joins, and sorting.

Conceptually:

\begin{lstlisting}[language=SQL]
CREATE INDEX idx_customer_id
ON orders(customer_id);
\end{lstlisting}

Indexes are particularly useful for frequently searched or joined columns, but
they have costs:

\begin{itemize}[leftmargin=*]
    \item additional storage;
    \item additional work during inserts and updates;
    \item not every query benefits from an index.
\end{itemize}

\textbf{Transactions}

Transactions group related modifications.

\begin{lstlisting}[language=SQL]
BEGIN;

UPDATE accounts
SET balance = balance - 100
WHERE account_id = 1;

UPDATE accounts
SET balance = balance + 100
WHERE account_id = 2;

COMMIT;
\end{lstlisting}

If something fails:

\begin{lstlisting}[language=SQL]
ROLLBACK;
\end{lstlisting}

Transactions are commonly described by \textbf{ACID}:

\begin{itemize}[leftmargin=*]
    \item \textbf{Atomicity}: the transaction succeeds as a unit or fails as a unit.
    \item \textbf{Consistency}: valid database rules remain satisfied.
    \item \textbf{Isolation}: concurrent transactions do not improperly interfere.
    \item \textbf{Durability}: committed changes persist.
\end{itemize}

\textbf{Practical tips}

\begin{itemize}[leftmargin=*]
    \item Always know what one row represents in every table.
    \item Check key uniqueness before joining.
    \item Compare row counts before and after important joins.
    \item Use CTEs to make complex transformations readable.
    \item Prefer explicit column names over \texttt{SELECT *}.
    \item Handle \texttt{NULL} intentionally rather than accidentally.
    \item Use \texttt{UNION ALL} when deduplication is unnecessary.
    \item Filter early when practical to avoid processing unnecessary rows.
    \item Verify the rows affected before running \texttt{UPDATE} or \texttt{DELETE}.
    \item Remember that date, string, and limit syntax varies somewhat by SQL dialect.
\end{itemize}

\textbf{Interview notes}

SQL interviews frequently test joins, aggregation, CTEs, window functions,
deduplication, top-$N$-per-group problems, and \texttt{NULL} behavior.

Be especially comfortable explaining:

\begin{itemize}[leftmargin=*]
    \item \texttt{WHERE} versus \texttt{HAVING};
    \item \texttt{INNER JOIN} versus \texttt{LEFT JOIN};
    \item \texttt{COUNT(*)} versus \texttt{COUNT(column)};
    \item \texttt{UNION} versus \texttt{UNION ALL};
    \item \texttt{ROW\_NUMBER} versus \texttt{RANK} versus \texttt{DENSE\_RANK};
    \item \texttt{GROUP BY} versus window functions;
    \item why joins can unexpectedly multiply rows;
    \item primary keys versus foreign keys;
    \item the basic purpose of indexes and transactions.
\end{itemize}


% ------------------------------------------------------------
% PYTHON INTEGRATION
% ------------------------------------------------------------

\subsubsection{Python Integration}

Python commonly interacts with relational databases through a database driver
or a database abstraction such as SQLAlchemy. For analytics, pandas is also
commonly used to move query results into DataFrames.

\textbf{Basic connection pattern}

Using SQLite:

\begin{lstlisting}[language=Python]
import sqlite3

conn = sqlite3.connect("database.db")

cursor = conn.cursor()

cursor.execute("""
    SELECT customer_id, revenue
    FROM customers
    WHERE revenue > ?
""", (1000,))

rows = cursor.fetchall()

cursor.close()
conn.close()
\end{lstlisting}

Common cursor operations:

\begin{lstlisting}[language=Python]
cursor.execute(query)

cursor.fetchone()
cursor.fetchmany(100)
cursor.fetchall()
\end{lstlisting}

\textbf{Context managers}

When supported, use context managers to manage connections safely:

\begin{lstlisting}[language=Python]
import sqlite3

with sqlite3.connect("database.db") as conn:
    rows = conn.execute(
        "SELECT * FROM customers"
    ).fetchall()
\end{lstlisting}

\textbf{Parameterized queries}

Do not create SQL queries using untrusted values through string interpolation.

Avoid:

\begin{lstlisting}[language=Python]
query = f"""
SELECT *
FROM users
WHERE username = '{username}'
"""
\end{lstlisting}

Instead use parameters:

\begin{lstlisting}[language=Python]
cursor.execute(
    """
    SELECT *
    FROM users
    WHERE username = ?
    """,
    (username,)
)
\end{lstlisting}

Parameterized queries correctly handle values and protect against SQL
injection.

Placeholder syntax differs across database drivers.

\textbf{pandas}

Read query results directly into a DataFrame:

\begin{lstlisting}[language=Python]
import pandas as pd

df = pd.read_sql_query(
    """
    SELECT
        customer_id,
        revenue
    FROM customers
    WHERE revenue > 1000
    """,
    conn
)
\end{lstlisting}

Write a DataFrame:

\begin{lstlisting}[language=Python]
df.to_sql(
    "results",
    conn,
    if_exists="append",
    index=False
)
\end{lstlisting}

Common \texttt{if\_exists} options:

\begin{center}
\begin{tabular}{ll}
\toprule
Option & Behavior \\
\midrule
\texttt{fail} & Raise an error if table exists \\
\texttt{replace} & Replace existing table \\
\texttt{append} & Add rows to existing table \\
\bottomrule
\end{tabular}
\end{center}

Be particularly careful with:

\begin{lstlisting}[language=Python]
if_exists="replace"
\end{lstlisting}

because replacing a table can destroy the existing table and may also lose
database-specific schema details.

\textbf{Push computation to SQL}

Avoid:

\begin{lstlisting}[language=Python]
df = pd.read_sql(
    "SELECT * FROM transactions",
    conn
)

df = df[df["year"] == 2026]
\end{lstlisting}

when the database can perform the filtering:

\begin{lstlisting}[language=Python]
df = pd.read_sql(
    """
    SELECT *
    FROM transactions
    WHERE year = 2026
    """,
    conn
)
\end{lstlisting}

A useful principle is:

\[
\boxed{\text{Move computation to the data before moving data to Python.}}
\]

This is especially important for large databases.

\textbf{Chunked reads}

For large query results:

\begin{lstlisting}[language=Python]
for chunk in pd.read_sql_query(
    query,
    conn,
    chunksize=10000
):
    process(chunk)
\end{lstlisting}

This prevents the entire result from being loaded into memory simultaneously.

\textbf{SQLAlchemy}

SQLAlchemy provides a common database interface.

\begin{lstlisting}[language=Python]
from sqlalchemy import create_engine

engine = create_engine(
    "sqlite:///database.db"
)
\end{lstlisting}

Use it with pandas:

\begin{lstlisting}[language=Python]
df = pd.read_sql(
    "SELECT * FROM customers",
    engine
)
\end{lstlisting}

Parameterized SQL:

\begin{lstlisting}[language=Python]
from sqlalchemy import text

query = text("""
    SELECT *
    FROM customers
    WHERE revenue > :min_revenue
""")

with engine.connect() as conn:
    result = conn.execute(
        query,
        {"min_revenue": 1000}
    )
\end{lstlisting}

\textbf{Transactions}

\begin{lstlisting}[language=Python]
with engine.begin() as conn:
    conn.execute(statement_1)
    conn.execute(statement_2)
\end{lstlisting}

The transaction commits when the block succeeds and rolls back if an exception
occurs.

\textbf{Practical tips}

\begin{itemize}[leftmargin=*]
    \item Use parameterized queries rather than string interpolation.
    \item Push filters, joins, and aggregation into SQL when practical.
    \item Do not load millions of unnecessary rows into pandas.
    \item Use chunking when results do not comfortably fit in memory.
    \item Close connections and cursors when they are no longer needed.
    \item Be cautious with \texttt{to\_sql(..., if\_exists="replace")}.
    \item Keep database credentials out of source code.
    \item Remember that pandas types and database types do not always map perfectly.
\end{itemize}

\textbf{Interview notes}

Understand the division of labor:

\[
\text{SQL}
=
\text{filter/join/aggregate near the data}
\]

\[
\text{Python}
=
\text{analysis/modeling/application logic}
\]

A common interview question is why filtering in SQL before loading into pandas
is preferable to retrieving an entire large table. The answer is reduced data
transfer, lower memory use, and use of the database engine's query optimizer.


% ------------------------------------------------------------
% R INTEGRATION
% ------------------------------------------------------------

\subsubsection{R Integration}

R commonly accesses relational databases through the \texttt{DBI} interface
and a database-specific backend.

The typical workflow is:

\[
\text{Connect}
\rightarrow
\text{Query}
\rightarrow
\text{Analyze}
\rightarrow
\text{Disconnect}
\]

\textbf{Basic DBI workflow}

Using SQLite:

\begin{lstlisting}[language=R]
library(DBI)
library(RSQLite)

con <- dbConnect(
    SQLite(),
    "database.db"
)

df <- dbGetQuery(
    con,
    "SELECT * FROM customers"
)

dbDisconnect(con)
\end{lstlisting}

\textbf{Useful DBI commands}

\begin{lstlisting}[language=R]
dbListTables(con)

dbListFields(
    con,
    "customers"
)

dbExistsTable(
    con,
    "customers"
)

dbReadTable(
    con,
    "customers"
)

dbGetQuery(
    con,
    "SELECT * FROM customers"
)

dbDisconnect(con)
\end{lstlisting}

\textbf{dbGetQuery versus dbExecute}

Use \texttt{dbGetQuery()} when SQL returns rows:

\begin{lstlisting}[language=R]
df <- dbGetQuery(
    con,
    "
    SELECT
        customer_id,
        revenue
    FROM customers
    "
)
\end{lstlisting}

Use \texttt{dbExecute()} for operations that modify data or schema:

\begin{lstlisting}[language=R]
dbExecute(
    con,
    "
    DELETE FROM customers
    WHERE inactive = 1
    "
)
\end{lstlisting}

\textbf{dbplyr}

\texttt{dbplyr} allows many familiar \texttt{dplyr} operations to be
translated into SQL.

\begin{lstlisting}[language=R]
library(dplyr)

customers <- tbl(
    con,
    "customers"
)

result <- customers |>
    filter(revenue > 1000) |>
    select(customer_id, revenue) |>
    arrange(desc(revenue))
\end{lstlisting}

A critical concept is that database-backed \texttt{dbplyr} operations are
typically \textbf{lazy}.

The result has not necessarily been loaded into R.

Inspect the generated SQL:

\begin{lstlisting}[language=R]
result |>
    show_query()
\end{lstlisting}

Load the result into R:

\begin{lstlisting}[language=R]
df <- result |>
    collect()
\end{lstlisting}

Therefore:

\[
\texttt{tbl()}
\rightarrow
\text{database-backed lazy query}
\]

while:

\[
\texttt{collect()}
\rightarrow
\text{materialized R object}
\]

\textbf{Push computation to the database}

Prefer:

\begin{lstlisting}[language=R]
result <- tbl(con, "transactions") |>
    filter(year == 2026) |>
    group_by(customer_id) |>
    summarize(
        total = sum(amount, na.rm = TRUE)
    ) |>
    collect()
\end{lstlisting}

over collecting the entire transaction table and then filtering it in R.

\textbf{Writing tables}

\begin{lstlisting}[language=R]
dbWriteTable(
    con,
    "results",
    df,
    overwrite = TRUE
)
\end{lstlisting}

Append instead:

\begin{lstlisting}[language=R]
dbWriteTable(
    con,
    "results",
    df,
    append = TRUE
)
\end{lstlisting}

Use \texttt{overwrite = TRUE} carefully because it replaces the existing
table.

\textbf{Practical tips}

\begin{itemize}[leftmargin=*]
    \item Use \texttt{show\_query()} when unsure how \texttt{dbplyr} translates an
    operation.
    \item Delay \texttt{collect()} until the data have been reduced to a manageable
    size.
    \item Push filters, joins, and summaries to the database.
    \item Disconnect from the database when finished.
    \item Avoid hard-coding database credentials in scripts.
    \item Be cautious when overwriting existing tables.
    \item Remember that not every R function has a SQL translation.
\end{itemize}

\textbf{Interview notes}

The most important concept is \textbf{lazy evaluation}. With
\texttt{dbplyr}, the R code often describes a query that executes in the
database rather than immediately creating an in-memory R data frame.

Be comfortable explaining why:

\begin{lstlisting}[language=R]
collect()
\end{lstlisting}

should usually occur \emph{after} filtering and aggregation rather than before.


% ============================================================
% DUCKDB
% ============================================================

\subsection{DuckDB}

\textbf{DuckDB} is an embedded analytical SQL database designed for fast
OLAP-style workloads. It is particularly useful for querying local analytical
data such as Parquet and CSV files without operating a separate database
server.

A useful analogy is:

\[
\text{SQLite}
:
\text{embedded transactional database}
\]

\[
\text{DuckDB}
:
\text{embedded analytical database}
\]

\textbf{Why use DuckDB?}

DuckDB is useful when:

\begin{itemize}[leftmargin=*]
    \item data are larger than is convenient for ordinary DataFrame workflows;
    \item data are stored in many CSV or Parquet files;
    \item SQL is convenient for the analysis;
    \item running a database server would be unnecessary overhead;
    \item only a subset of a large dataset needs to be queried.
\end{itemize}

\textbf{CLI}

Start an in-memory session:

\begin{lstlisting}
duckdb
\end{lstlisting}

Open or create a persistent database:

\begin{lstlisting}
duckdb analytics.duckdb
\end{lstlisting}

Useful commands:

\begin{lstlisting}
.tables
.schema
.schema customers
.describe customers
.help
.quit
\end{lstlisting}

\textbf{Query CSV directly}

\begin{lstlisting}[language=SQL]
SELECT *
FROM 'customers.csv'
LIMIT 10;
\end{lstlisting}

Explicit CSV reader:

\begin{lstlisting}[language=SQL]
SELECT *
FROM read_csv_auto('customers.csv');
\end{lstlisting}

Create a persistent table:

\begin{lstlisting}[language=SQL]
CREATE TABLE customers AS
SELECT *
FROM read_csv_auto('customers.csv');
\end{lstlisting}

\textbf{Query Parquet directly}

\begin{lstlisting}[language=SQL]
SELECT *
FROM 'transactions.parquet'
LIMIT 10;
\end{lstlisting}

or:

\begin{lstlisting}[language=SQL]
SELECT *
FROM read_parquet('transactions.parquet');
\end{lstlisting}

Query many files:

\begin{lstlisting}[language=SQL]
SELECT *
FROM read_parquet('transactions/*.parquet');
\end{lstlisting}

Aggregate without importing the files first:

\begin{lstlisting}[language=SQL]
SELECT
    customer_id,
    SUM(amount) AS total_spend
FROM read_parquet('transactions/*.parquet')
WHERE transaction_date >= '2026-01-01'
GROUP BY customer_id;
\end{lstlisting}

This is one of DuckDB's most useful everyday capabilities.

\textbf{Views over files}

Create a reusable view:

\begin{lstlisting}[language=SQL]
CREATE VIEW transactions AS
SELECT *
FROM read_parquet(
    'transactions/*.parquet'
);
\end{lstlisting}

Then:

\begin{lstlisting}[language=SQL]
SELECT
    customer_id,
    SUM(amount) AS total_spend
FROM transactions
GROUP BY customer_id;
\end{lstlisting}

\textbf{Export to Parquet}

\begin{lstlisting}[language=SQL]
COPY (
    SELECT *
    FROM transactions
    WHERE transaction_date >= '2026-01-01'
)
TO 'transactions_2026.parquet'
(FORMAT PARQUET);
\end{lstlisting}

Export CSV:

\begin{lstlisting}[language=SQL]
COPY (
    SELECT *
    FROM transactions
)
TO 'transactions.csv'
(HEADER, DELIMITER ',');
\end{lstlisting}

\textbf{Parquet}

Parquet is a \textbf{columnar} file format.

For an analytical query such as:

\begin{lstlisting}[language=SQL]
SELECT
    customer_id,
    amount
FROM read_parquet('large.parquet')
WHERE year = 2026;
\end{lstlisting}

DuckDB can often avoid reading irrelevant data.

Two important concepts are:

\begin{itemize}[leftmargin=*]
    \item \textbf{Projection pushdown}: read only required columns.
    \item \textbf{Predicate pushdown}: avoid reading data that cannot satisfy filters.
\end{itemize}

This is why DuckDB + Parquet can be substantially more efficient than loading
an entire CSV into memory.

\textbf{EXPLAIN}

Inspect the query plan:

\begin{lstlisting}[language=SQL]
EXPLAIN
SELECT
    customer_id,
    SUM(amount)
FROM transactions
GROUP BY customer_id;
\end{lstlisting}

Inspect execution information:

\begin{lstlisting}[language=SQL]
EXPLAIN ANALYZE
SELECT
    customer_id,
    SUM(amount)
FROM transactions
GROUP BY customer_id;
\end{lstlisting}

\textbf{DuckDB versus pandas}

A useful rule of thumb is:

\[
\text{pandas}
\rightarrow
\text{in-memory manipulation}
\]

\[
\text{DuckDB}
\rightarrow
\text{SQL analytical engine}
\]

DuckDB is particularly attractive when only a small subset or aggregation of
a large file collection is needed.

\textbf{DuckDB versus server databases}

\begin{center}
\begin{tabularx}{\textwidth}{lXX}
\toprule
 & DuckDB & Server Database \\
\midrule
Architecture & Embedded & Client/server \\
Setup & Minimal & Server infrastructure \\
Primary strength & Local analytics & Shared persistent workloads \\
Parquet querying & Excellent & System dependent \\
Multi-user workloads & Not primary focus & Common use case \\
\bottomrule
\end{tabularx}
\end{center}

\textbf{Practical tips}

\begin{itemize}[leftmargin=*]
    \item Prefer Parquet over CSV for repeated analytical work when possible.
    \item Query files directly before deciding whether they need to become tables.
    \item Use wildcards to query partitioned file collections.
    \item Select only the columns required by the analysis.
    \item Filter early to take advantage of predicate pushdown.
    \item Use \texttt{EXPLAIN ANALYZE} when investigating slow queries.
    \item Do not deploy DuckDB as though it were a traditional high-concurrency
    application database.
\end{itemize}

\textbf{Interview notes}

The core distinction is:

\[
\boxed{\text{DuckDB is an embedded OLAP database.}}
\]

Be able to explain why columnar storage is useful for analytics, why Parquet
pairs well with DuckDB, and what predicate and projection pushdown mean.

A useful comparison is that PostgreSQL/MySQL commonly operate as database
servers, SQLite is an embedded database commonly associated with transactional
application workloads, and DuckDB is optimized for embedded analytical
workloads.


% ============================================================
% OMOP COMMON DATA MODEL
% ============================================================

\subsection{OMOP Common Data Model}

\textbf{OMOP (Observational Medical Outcomes Partnership) Common Data Model}
is a standardized relational data model maintained by the OHDSI community for
observational healthcare data.

Its purpose is to transform heterogeneous healthcare data into a common
structure and vocabulary so that analyses can be reused across datasets and
institutions.

A useful mental model is:

\[
\text{EHR / Claims / Registry}
\rightarrow
\text{ETL}
\rightarrow
\text{OMOP CDM}
\rightarrow
\text{Standardized Analysis}
\]

\textbf{FHIR versus OMOP}

The distinction is important:

\[
\boxed{\text{FHIR} \rightarrow \text{interoperability}}
\]

\[
\boxed{\text{OMOP} \rightarrow \text{observational analytics}}
\]

FHIR is resource-oriented and primarily designed for healthcare information
exchange. OMOP is relational and primarily designed for standardized
observational research and analytics.

\textbf{Core clinical tables}

\begin{center}
\begin{tabularx}{\textwidth}{lX}
\toprule
Table & Purpose \\
\midrule
PERSON & Patient/person demographics \\
VISIT\_OCCURRENCE & Healthcare encounters \\
CONDITION\_OCCURRENCE & Diagnoses and conditions \\
DRUG\_EXPOSURE & Medication exposures \\
PROCEDURE\_OCCURRENCE & Procedures \\
MEASUREMENT & Labs and quantitative measurements \\
OBSERVATION & Clinical facts not represented elsewhere \\
DEVICE\_EXPOSURE & Medical device exposure \\
SPECIMEN & Biological specimens \\
DEATH & Mortality information \\
\bottomrule
\end{tabularx}
\end{center}

A simplified relationship is:

\[
\texttt{PERSON}
\rightarrow
\begin{cases}
\texttt{VISIT\_OCCURRENCE} \\
\texttt{CONDITION\_OCCURRENCE} \\
\texttt{DRUG\_EXPOSURE} \\
\texttt{PROCEDURE\_OCCURRENCE} \\
\texttt{MEASUREMENT}
\end{cases}
\]

Most clinical event tables therefore contain a \texttt{person\_id}.

\textbf{PERSON}

The \texttt{PERSON} table contains one row per person.

Common concepts include:

\begin{lstlisting}
person_id
gender_concept_id
year_of_birth
month_of_birth
day_of_birth
race_concept_id
ethnicity_concept_id
\end{lstlisting}

The exact availability of demographic fields depends on the source data.

\textbf{VISIT\_OCCURRENCE}

Visits represent healthcare encounters.

Conceptually:

\begin{lstlisting}
person_id
visit_concept_id
visit_start_date
visit_end_date
\end{lstlisting}

Clinical events may reference the visit during which they occurred through
\texttt{visit\_occurrence\_id}.

\textbf{CONDITION\_OCCURRENCE}

Conditions and diagnoses are represented as events:

\begin{lstlisting}[language=SQL]
SELECT
    person_id,
    condition_concept_id,
    condition_start_date
FROM condition_occurrence
WHERE person_id = 123;
\end{lstlisting}

\textbf{MEASUREMENT}

Measurements commonly represent laboratory results and other quantitative
clinical measurements.

Important concepts include:

\begin{lstlisting}
person_id
measurement_concept_id
measurement_date
value_as_number
value_as_concept_id
unit_concept_id
\end{lstlisting}

When analyzing measurements, do not ignore the unit.

For example:

\[
100\ \text{mg/dL}
\neq
100\ \text{mmol/L}
\]

\textbf{Standard concepts}

OMOP relies heavily on standardized concept identifiers.

Conceptually:

\[
\text{Source Code}
\rightarrow
\text{Standard Concept}
\]

Clinical tables often contain both standardized and source representations.

For example:

\begin{lstlisting}
condition_concept_id
condition_source_concept_id
condition_source_value
\end{lstlisting}

The standardized concept enables portable analysis, while source fields
preserve information from the original system.

\textbf{CONCEPT}

The \texttt{CONCEPT} table provides information about vocabulary concepts.

Useful fields include:

\begin{lstlisting}
concept_id
concept_name
domain_id
vocabulary_id
concept_class_id
standard_concept
concept_code
\end{lstlisting}

Join a condition to its human-readable concept:

\begin{lstlisting}[language=SQL]
SELECT
    c.person_id,
    c.condition_start_date,
    concept.concept_name
FROM condition_occurrence AS c
JOIN concept
    ON c.condition_concept_id =
       concept.concept_id
WHERE c.person_id = 123;
\end{lstlisting}

\textbf{Vocabulary tables}

Important vocabulary-related tables include:

\begin{center}
\begin{tabularx}{\textwidth}{lX}
\toprule
Table & Purpose \\
\midrule
CONCEPT & Individual standardized concepts \\
VOCABULARY & Vocabulary definitions \\
DOMAIN & Clinical concept domains \\
CONCEPT\_RELATIONSHIP & Relationships among concepts \\
CONCEPT\_ANCESTOR & Hierarchical ancestor/descendant relationships \\
CONCEPT\_SYNONYM & Alternative concept names \\
\bottomrule
\end{tabularx}
\end{center}

\textbf{Concept hierarchies}

Suppose an analysis needs patients with any diagnosis belonging to a broader
disease family.

Rather than manually listing every descendant code,
\texttt{CONCEPT\_ANCESTOR} can represent the hierarchy:

\[
\text{Broad Ancestor}
\rightarrow
\text{Specific Descendants}
\]

For example:

\begin{lstlisting}[language=SQL]
SELECT
    descendant_concept_id
FROM concept_ancestor
WHERE ancestor_concept_id = 123456;
\end{lstlisting}

Use this to identify conditions belonging to the broader concept:

\begin{lstlisting}[language=SQL]
SELECT DISTINCT
    c.person_id
FROM condition_occurrence AS c
JOIN concept_ancestor AS ca
    ON c.condition_concept_id =
       ca.descendant_concept_id
WHERE ca.ancestor_concept_id = 123456;
\end{lstlisting}

\textbf{Domains}

Standard concepts belong to domains such as:

\begin{itemize}[leftmargin=*]
    \item Condition;
    \item Drug;
    \item Procedure;
    \item Measurement;
    \item Observation;
    \item Device.
\end{itemize}

The domain helps determine where a standardized clinical concept belongs in
the OMOP structure.

\textbf{Cohort construction}

A common OMOP task is defining a patient cohort.

For example:

\begin{lstlisting}[language=SQL]
SELECT DISTINCT
    person_id
FROM condition_occurrence
WHERE condition_concept_id IN (
    12345,
    67890
);
\end{lstlisting}

Add a temporal criterion:

\begin{lstlisting}[language=SQL]
SELECT DISTINCT
    person_id
FROM condition_occurrence
WHERE condition_concept_id = 12345
  AND condition_start_date >= '2026-01-01';
\end{lstlisting}

Real cohort definitions may combine:

\begin{itemize}[leftmargin=*]
    \item diagnoses;
    \item medications;
    \item measurements;
    \item procedures;
    \item demographics;
    \item temporal relationships;
    \item inclusion/exclusion windows.
\end{itemize}

\textbf{Grain matters}

Different OMOP tables have different grains.

For example:

\[
\texttt{PERSON}
=
1\text{ row/person}
\]

while:

\[
\texttt{CONDITION\_OCCURRENCE}
=
1\text{ row/condition occurrence}
\]

and:

\[
\texttt{MEASUREMENT}
=
1\text{ row/measurement}
\]

Joining these directly can produce many rows per patient.

Always ask:

\begin{quote}
What does one row represent in this OMOP table?
\end{quote}

\textbf{Standardization does not imply perfect data}

An OMOP database can be structurally valid while still containing problems
with:

\begin{itemize}[leftmargin=*]
    \item missing data;
    \item incomplete vocabulary mappings;
    \item implausible values;
    \item duplicate events;
    \item inconsistent units;
    \item source-specific coding practices;
    \item incomplete capture of healthcare utilization.
\end{itemize}

Therefore:

\[
\boxed{\text{Common schema} \neq \text{guaranteed data quality}}
\]

\textbf{Practical tips}

\begin{itemize}[leftmargin=*]
    \item Start by identifying the clinical domain and corresponding OMOP table.
    \item Use standardized concept IDs for portable analyses.
    \item Preserve source concepts when investigating unexpected mappings.
    \item Join to \texttt{CONCEPT} when you need readable concept names.
    \item Use \texttt{CONCEPT\_ANCESTOR} instead of manually enumerating large
    clinical concept families when appropriate.
    \item Check units before analyzing measurements.
    \item Determine table grain before joining multiple clinical domains.
    \item Use \texttt{DISTINCT} thoughtfully; it can hide an incorrect many-to-many
    join rather than actually fixing it.
    \item Validate cohort logic against clinical expectations, not just SQL syntax.
\end{itemize}

\textbf{Interview notes}

Be able to explain why OMOP exists: different healthcare organizations use
different source systems and coding conventions, so OMOP maps those data into
a common relational model and standardized vocabulary.

The most useful distinctions to remember are:

\begin{itemize}[leftmargin=*]
    \item OMOP versus FHIR;
    \item source concepts versus standard concepts;
    \item person-level versus event-level tables;
    \item \texttt{CONCEPT} versus \texttt{CONCEPT\_ANCESTOR};
    \item structural standardization versus actual data quality.
\end{itemize}

For analytical work, the combination of \textbf{SQL + OMOP vocabulary
knowledge + careful cohort construction} is generally more important than
memorizing every OMOP table.

% ============================================================
% 3. ANALYSIS AND MODELING
% ============================================================

\section{Analysis and Modeling}

% ============================================================
% GOOGLE COLAB
% ============================================================

\subsection{Google Colab}

\textbf{Google Colab} is a hosted Jupyter notebook environment useful for
exploratory analysis, prototyping, teaching, and temporary GPU/TPU compute.

\textbf{Common shortcuts}

\begin{center}
\begin{tabularx}{\textwidth}{lX}
\toprule
Shortcut & Action \\
\midrule
Shift + Enter & Run cell and move to next cell \\
Ctrl/Cmd + Enter & Run current cell \\
Alt/Option + Enter & Run cell and insert a new cell \\
Ctrl/Cmd + M, B & Insert cell below \\
Ctrl/Cmd + M, A & Insert cell above \\
Ctrl/Cmd + M, D & Delete current cell \\
\bottomrule
\end{tabularx}
\end{center}

\textbf{Shell commands}

Prefix shell commands with \texttt{!}:

\begin{lstlisting}
!pwd
!ls
!pip install package_name
!git clone https://github.com/user/project.git
\end{lstlisting}

The default working directory is usually:

\begin{lstlisting}
/content
\end{lstlisting}

\textbf{Useful magic commands}

\begin{lstlisting}
%pwd
%cd /content/project
%time expression
%timeit expression
%who
%whos
\end{lstlisting}

Time an entire cell:

\begin{lstlisting}
%%time

# code
\end{lstlisting}

\textbf{Google Drive}

The runtime filesystem is temporary. Mount Drive when persistent storage is
needed:

\begin{lstlisting}[language=Python]
from google.colab import drive

drive.mount("/content/drive")
\end{lstlisting}

Files are commonly available under:

\begin{lstlisting}
/content/drive/MyDrive/
\end{lstlisting}

\textbf{File upload/download}

\begin{lstlisting}[language=Python]
from google.colab import files

uploaded = files.upload()

files.download("results.csv")
\end{lstlisting}

\textbf{Packages}

\begin{lstlisting}
!pip install transformers
!pip install -q datasets accelerate
\end{lstlisting}

For reproducible work, pin important versions:

\begin{lstlisting}
!pip install package_name==1.2.3
\end{lstlisting}

A restarted runtime may require packages to be installed again.

\textbf{GPU}

Check the assigned GPU:

\begin{lstlisting}
!nvidia-smi
\end{lstlisting}

Check availability with PyTorch:

\begin{lstlisting}[language=Python]
import torch

torch.cuda.is_available()
torch.cuda.get_device_name(0)
\end{lstlisting}

Move a model to the available device:

\begin{lstlisting}[language=Python]
device = torch.device(
    "cuda" if torch.cuda.is_available() else "cpu"
)

model = model.to(device)
\end{lstlisting}

Selecting a GPU runtime does not automatically make all computations use the
GPU.

\textbf{Good notebook structure}

\begin{lstlisting}
1. Imports
2. Configuration and seeds
3. Data loading
4. Data processing
5. Modeling
6. Evaluation
7. Results / export
\end{lstlisting}

Before sharing a notebook, test:

\begin{lstlisting}
Restart runtime -> Run all
\end{lstlisting}

This catches hidden dependencies caused by running cells out of order.

\textbf{Practical tips}

\begin{itemize}[leftmargin=*]
    \item Assume files in \texttt{/content} are temporary.
    \item Save important results and checkpoints to persistent storage.
    \item Keep installation and configuration cells near the top.
    \item Restart and run all cells before sharing a notebook.
    \item Avoid relying on variables created by out-of-order cell execution.
    \item Use \texttt{!nvidia-smi} to verify accelerator availability.
    \item Do not hard-code secrets into notebooks.
    \item Save checkpoints during long training jobs.
\end{itemize}

\textbf{Interview notes}

Colab is a hosted development environment, not a production deployment
platform. Be comfortable discussing ephemeral runtimes, reproducibility,
notebook state, package environments, and CPU/GPU execution.


% ============================================================
% SAS
% ============================================================

\subsection{SAS}

\textbf{SAS} is a statistical programming and data-management environment
widely used in healthcare, clinical research, government, finance, and
regulated enterprise settings.

The two fundamental building blocks are:

\begin{itemize}[leftmargin=*]
    \item \textbf{DATA step}: create, transform, and filter data.
    \item \textbf{PROC step}: analyze, summarize, model, sort, or report data.
\end{itemize}

\textbf{DATA step}

\begin{lstlisting}
DATA adults;
    SET patients;
    WHERE age >= 18;
RUN;
\end{lstlisting}

Create variables:

\begin{lstlisting}
DATA patients2;
    SET patients;

    bmi = weight_kg / (height_m ** 2);

    IF bmi >= 30 THEN obese = 1;
    ELSE obese = 0;
RUN;
\end{lstlisting}

\textbf{Libraries}

Assign a permanent library:

\begin{lstlisting}
LIBNAME study "/path/to/data";
\end{lstlisting}

Reference a dataset:

\begin{lstlisting}
study.patients
\end{lstlisting}

Temporary datasets are commonly stored in:

\begin{lstlisting}
WORK
\end{lstlisting}

The \texttt{WORK} library disappears when the SAS session ends.

\textbf{Inspect data}

\begin{lstlisting}
PROC CONTENTS DATA=study.patients;
RUN;

PROC PRINT DATA=study.patients (OBS=10);
RUN;
\end{lstlisting}

\textbf{KEEP and DROP}

\begin{lstlisting}
DATA analysis;
    SET patients;
    KEEP patient_id age treatment outcome;
RUN;
\end{lstlisting}

or:

\begin{lstlisting}
DATA analysis;
    SET patients;
    DROP address phone_number;
RUN;
\end{lstlisting}

\textbf{WHERE versus IF}

\begin{lstlisting}
WHERE age >= 18;
\end{lstlisting}

filters observations as they are read and generally uses variables already
present in the input data.

\begin{lstlisting}
IF age >= 18;
\end{lstlisting}

operates during DATA-step processing and can use variables created within the
step.

\textbf{Sorting}

\begin{lstlisting}
PROC SORT DATA=patients
          OUT=patients_sorted;
    BY patient_id visit_date;
RUN;
\end{lstlisting}

\textbf{BY-group processing}

After sorting:

\begin{lstlisting}
DATA first_visit;
    SET patients_sorted;

    BY patient_id;

    IF FIRST.patient_id;
RUN;
\end{lstlisting}

Useful temporary indicators:

\begin{lstlisting}
FIRST.patient_id
LAST.patient_id
\end{lstlisting}

Keep the latest record:

\begin{lstlisting}
PROC SORT DATA=patients;
    BY patient_id visit_date;
RUN;

DATA latest;
    SET patients;
    BY patient_id;

    IF LAST.patient_id;
RUN;
\end{lstlisting}

\textbf{Missing values}

Numeric missing value:

\begin{lstlisting}
.
\end{lstlisting}

Check missingness:

\begin{lstlisting}
IF MISSING(age) THEN ...
\end{lstlisting}

Character missing values are represented by blanks.

\textbf{PROC FREQ}

\begin{lstlisting}
PROC FREQ DATA=patients;
    TABLES treatment;
RUN;
\end{lstlisting}

Cross-tabulation:

\begin{lstlisting}
PROC FREQ DATA=patients;
    TABLES treatment*outcome / CHISQ;
RUN;
\end{lstlisting}

\textbf{PROC MEANS}

\begin{lstlisting}
PROC MEANS DATA=patients
           N MEAN STD MEDIAN MIN MAX;
    VAR age bmi;
RUN;
\end{lstlisting}

Grouped summary:

\begin{lstlisting}
PROC MEANS DATA=patients;
    CLASS treatment;
    VAR outcome;
RUN;
\end{lstlisting}

\textbf{PROC SQL}

\begin{lstlisting}
PROC SQL;
    SELECT
        treatment,
        COUNT(*) AS n,
        AVG(outcome) AS mean_outcome
    FROM patients
    GROUP BY treatment;
QUIT;
\end{lstlisting}

Create a dataset:

\begin{lstlisting}
PROC SQL;
    CREATE TABLE adults AS
    SELECT *
    FROM patients
    WHERE age >= 18;
QUIT;
\end{lstlisting}

\textbf{Linear regression}

\begin{lstlisting}
PROC REG DATA=patients;
    MODEL outcome = age bmi treatment;
RUN;
QUIT;
\end{lstlisting}

\textbf{Logistic regression}

\begin{lstlisting}
PROC LOGISTIC DATA=patients;
    CLASS treatment / PARAM=REF;

    MODEL outcome(EVENT='1') =
        age bmi treatment;
RUN;
\end{lstlisting}

\textbf{Formats}

Formats change how values are displayed without changing the underlying stored
values.

\begin{lstlisting}
PROC FORMAT;
    VALUE treatment_fmt
        0 = "Control"
        1 = "Treatment";
RUN;

DATA patients;
    SET patients;
    FORMAT treatment treatment_fmt.;
RUN;
\end{lstlisting}

\textbf{Macro variables}

\begin{lstlisting}
%LET year = 2026;

PROC PRINT DATA=patients;
    WHERE year = &year;
RUN;
\end{lstlisting}

Simple macro:

\begin{lstlisting}
%MACRO summarize(var);

PROC MEANS DATA=patients;
    VAR &var;
RUN;

%MEND;

%summarize(age)
\end{lstlisting}

Macros perform text/code generation and are different from DATA-step
variables.

\textbf{Practical tips}

\begin{itemize}[leftmargin=*]
    \item Know whether data are temporary (\texttt{WORK}) or permanent.
    \item Use \texttt{PROC CONTENTS} when opening an unfamiliar dataset.
    \item Sort before relying on BY-group processing.
    \item Know \texttt{FIRST.} and \texttt{LAST.} for grouped records.
    \item Distinguish displayed formats from stored values.
    \item Check the SAS log for warnings and errors after execution.
    \item Handle missing values explicitly.
\end{itemize}

\textbf{Interview notes}

Know DATA versus PROC steps, \texttt{SET}, \texttt{WHERE} versus
\texttt{IF}, libraries, BY-group processing, \texttt{FIRST.}/\texttt{LAST.},
formats, missing values, \texttt{PROC SQL}, \texttt{PROC FREQ}, and
\texttt{PROC MEANS}.


% ============================================================
% STATA
% ============================================================

\subsection{Stata}

\textbf{Stata} is a statistical programming environment widely used in
economics, econometrics, epidemiology, public policy, and social science.

General command structure:

\begin{lstlisting}
command varlist [if] [in], options
\end{lstlisting}

Example:

\begin{lstlisting}
summarize income if age >= 18, detail
\end{lstlisting}

\textbf{Load and save}

\begin{lstlisting}
use "data.dta", clear

save "clean_data.dta", replace

import delimited "data.csv", clear

export delimited using "results.csv", replace
\end{lstlisting}

\textbf{Inspect data}

\begin{lstlisting}
describe
codebook
summarize
summarize income, detail
list in 1/10
browse
count
tabulate treatment
\end{lstlisting}

Cross-tabulation:

\begin{lstlisting}
tabulate treatment outcome
\end{lstlisting}

\textbf{Generate and replace}

\begin{lstlisting}
generate bmi = weight_kg / (height_m^2)

generate adult = age >= 18

replace income = . if income < 0
\end{lstlisting}

\texttt{gen} is a common abbreviation for \texttt{generate}.

\textbf{Missing values}

Standard numeric missing:

\begin{lstlisting}
.
\end{lstlisting}

Extended missing values:

\begin{lstlisting}
.a .b .c ... .z
\end{lstlisting}

Important Stata behavior: numeric missing values compare greater than ordinary
numeric values.

Therefore:

\begin{lstlisting}
keep if age > 18
\end{lstlisting}

may retain missing ages.

Safer:

\begin{lstlisting}
keep if age > 18 & !missing(age)
\end{lstlisting}

\textbf{Keep and drop}

Variables:

\begin{lstlisting}
keep patient_id age treatment outcome

drop address phone_number
\end{lstlisting}

Observations:

\begin{lstlisting}
keep if age >= 18

drop if missing(outcome)
\end{lstlisting}

\textbf{Sort and bysort}

\begin{lstlisting}
sort patient_id date

bysort patient_id (date): gen visit_number = _n
\end{lstlisting}

Within a group:

\begin{itemize}[leftmargin=*]
    \item \texttt{\_n}: current observation number.
    \item \texttt{\_N}: total observations in the group.
\end{itemize}

Keep first record per patient:

\begin{lstlisting}
bysort patient_id (date): keep if _n == 1
\end{lstlisting}

Count observations per patient:

\begin{lstlisting}
bysort patient_id: gen n_visits = _N
\end{lstlisting}

\textbf{egen}

Grouped mean:

\begin{lstlisting}
bysort state: egen mean_income = mean(income)
\end{lstlisting}

Group identifier:

\begin{lstlisting}
egen firm_id = group(firm_name)
\end{lstlisting}

Row total:

\begin{lstlisting}
egen total_score = rowtotal(score1 score2 score3)
\end{lstlisting}

\textbf{Collapse}

\begin{lstlisting}
collapse (mean) income age ///
         (count) n=id, by(state)
\end{lstlisting}

\texttt{collapse} replaces the current dataset in memory.

Use:

\begin{lstlisting}
preserve

collapse (mean) income, by(state)

restore
\end{lstlisting}

when the original data are still needed.

\textbf{Reshape}

Wide to long:

\begin{lstlisting}
reshape long income, i(id) j(year)
\end{lstlisting}

Long to wide:

\begin{lstlisting}
reshape wide income, i(id) j(year)
\end{lstlisting}

\textbf{Merge}

\begin{lstlisting}
merge 1:1 patient_id using "demographics.dta"

merge 1:m id using "visits.dta"

merge m:1 firm_id using "firm_data.dta"
\end{lstlisting}

Always inspect:

\begin{lstlisting}
tab _merge
\end{lstlisting}

Typical values:

\begin{center}
\begin{tabular}{cl}
\toprule
\texttt{\_merge} & Meaning \\
\midrule
1 & Master only \\
2 & Using only \\
3 & Matched \\
\bottomrule
\end{tabular}
\end{center}

Avoid \texttt{merge m:m} for ordinary relational data. It usually indicates
that the intended key structure has not been resolved correctly.

\textbf{Linear regression}

\begin{lstlisting}
regress y x1 x2 x3

regress y x1 x2 x3, vce(robust)

regress y x1 x2 x3, vce(cluster firm_id)
\end{lstlisting}

\textbf{Logistic regression}

\begin{lstlisting}
logit y x1 x2 x3

logistic y x1 x2 x3
\end{lstlisting}

\textbf{Factor variables}

\begin{lstlisting}
regress y x1 i.industry

regress y c.age i.treatment

regress y c.age##i.treatment
\end{lstlisting}

Useful distinction:

\begin{itemize}[leftmargin=*]
    \item \texttt{\#}: interaction term.
    \item \texttt{\#\#}: main effects plus interaction.
\end{itemize}

\textbf{Panel data}

\begin{lstlisting}
xtset firm_id year

xtreg y x1 x2, fe

xtreg y x1 x2, re
\end{lstlisting}

\textbf{Time series}

\begin{lstlisting}
tsset date

generate lag_y = L.y

generate diff_y = D.y

generate lead_y = F.y
\end{lstlisting}

\textbf{Local macros}

\begin{lstlisting}
local controls age income education

regress y treatment `controls'
\end{lstlisting}

Prefer local macros over global macros when possible.

\textbf{Do-files}

A reproducible Stata workflow should usually live in a \texttt{.do} file:

\begin{lstlisting}
clear all
set more off

cd "/path/to/project"

log using "analysis.log", replace

* Load data
* Clean data
* Analyze
* Export results

log close
\end{lstlisting}

\textbf{Practical tips}

\begin{itemize}[leftmargin=*]
    \item Work from do-files rather than relying on command history.
    \item Inspect \texttt{\_merge} after every merge.
    \item Avoid \texttt{merge m:m} unless there is a specific reason.
    \item Remember Stata's unusual numeric missing-value ordering.
    \item Use \texttt{preserve}/\texttt{restore} around destructive temporary
    operations.
    \item Use factor-variable notation rather than manually creating every dummy.
    \item Declare panel or time-series structure before using corresponding
    operators.
\end{itemize}

\textbf{Interview notes}

Know \texttt{generate}, \texttt{replace}, \texttt{egen}, \texttt{bysort},
\texttt{\_n}, \texttt{\_N}, \texttt{merge}, \texttt{reshape}, and
\texttt{collapse}.

For econometrics roles, also know factor variables, robust versus clustered
standard errors, panel declarations, fixed versus random effects, and basic
lag/difference operators.


% ============================================================
% JAX
% ============================================================

\subsection{JAX}

\textbf{JAX} is a numerical computing library with a NumPy-like interface,
automatic differentiation, JIT compilation, vectorization, and accelerator
support.

The three transformations to remember are:

\begin{itemize}[leftmargin=*]
    \item \texttt{grad}: automatic differentiation.
    \item \texttt{jit}: just-in-time compilation.
    \item \texttt{vmap}: automatic vectorization.
\end{itemize}

\textbf{Arrays}

\begin{lstlisting}[language=Python]
import jax
import jax.numpy as jnp

x = jnp.array([1.0, 2.0, 3.0])

jnp.mean(x)
jnp.sum(x)
jnp.exp(x)
jnp.dot(x, x)
\end{lstlisting}

JAX intentionally resembles NumPy:

\begin{lstlisting}[language=Python]
import numpy as np
import jax.numpy as jnp
\end{lstlisting}

but the execution models differ.

\textbf{Automatic differentiation}

\begin{lstlisting}[language=Python]
from jax import grad

def f(x):
    return x**2 + 3*x

df = grad(f)

df(2.0)
\end{lstlisting}

For a loss:

\begin{lstlisting}[language=Python]
def loss(w, x, y):
    pred = x @ w
    return jnp.mean((pred - y)**2)

grad_loss = grad(loss)

grads = grad_loss(w, x, y)
\end{lstlisting}

By default, \texttt{grad()} differentiates with respect to the first argument.

Specify another argument with:

\begin{lstlisting}[language=Python]
grad(loss, argnums=0)
\end{lstlisting}

\textbf{Value and gradient together}

\begin{lstlisting}[language=Python]
from jax import value_and_grad

loss_and_grad = value_and_grad(loss)

value, grads = loss_and_grad(w, x, y)
\end{lstlisting}

\textbf{JIT compilation}

\begin{lstlisting}[language=Python]
from jax import jit

@jit
def f(x):
    return jnp.sin(x) + jnp.cos(x)
\end{lstlisting}

or:

\begin{lstlisting}[language=Python]
compiled_f = jit(f)
\end{lstlisting}

The first compatible call includes tracing and compilation overhead. Later
calls can be much faster.

Do not benchmark only the first JIT call when measuring steady-state
performance.

\textbf{Vectorization}

Suppose a function handles one observation:

\begin{lstlisting}[language=Python]
def predict(w, x):
    return x @ w
\end{lstlisting}

Vectorize it:

\begin{lstlisting}[language=Python]
from jax import vmap

batch_predict = vmap(
    predict,
    in_axes=(None, 0)
)

predictions = batch_predict(w, X)
\end{lstlisting}

Here:

\begin{itemize}[leftmargin=*]
    \item \texttt{None}: reuse the same \texttt{w}.
    \item \texttt{0}: vectorize across axis 0 of \texttt{X}.
\end{itemize}

\textbf{Immutable arrays}

Do not use NumPy-style mutation:

\begin{lstlisting}[language=Python]
# Avoid
x[0] = 10
\end{lstlisting}

Use functional updates:

\begin{lstlisting}[language=Python]
x = x.at[0].set(10)

x = x.at[0].add(5)

x = x.at[idx].set(values)
\end{lstlisting}

\textbf{Randomness}

JAX uses explicit pseudo-random keys:

\begin{lstlisting}[language=Python]
from jax import random

key = random.key(42)

x = random.normal(
    key,
    shape=(100,)
)
\end{lstlisting}

Split keys when independent randomness is required:

\begin{lstlisting}[language=Python]
key, subkey = random.split(key)

x = random.normal(
    subkey,
    shape=(100,)
)
\end{lstlisting}

Do not repeatedly reuse the same key when independent draws are intended.

\textbf{PyTrees}

Model parameters can be stored in nested structures:

\begin{lstlisting}[language=Python]
params = {
    "layer1": {
        "weights": jnp.ones((10, 5)),
        "bias": jnp.zeros(5)
    },
    "layer2": {
        "weights": jnp.ones((5, 1)),
        "bias": jnp.zeros(1)
    }
}
\end{lstlisting}

JAX transformations understand many nested structures such as dictionaries,
lists, and tuples as PyTrees.

\textbf{Devices}

\begin{lstlisting}[language=Python]
jax.devices()

jax.default_backend()
\end{lstlisting}

JAX can execute on CPU, GPU, or TPU depending on the environment.

\textbf{JAX versus NumPy}

\begin{center}
\begin{tabularx}{\textwidth}{lXX}
\toprule
Feature & NumPy & JAX \\
\midrule
Array interface & Standard & NumPy-like \\
Autodiff & No & Yes \\
JIT compilation & No & Yes \\
GPU/TPU & Not native & Supported \\
Mutation & Supported & Functional style \\
Randomness & Stateful-style & Explicit keys \\
\bottomrule
\end{tabularx}
\end{center}

\textbf{Practical tips}

\begin{itemize}[leftmargin=*]
    \item Write pure, array-oriented functions when possible.
    \item Do not mutate JAX arrays in place.
    \item Split PRNG keys deliberately.
    \item Expect compilation overhead on the first JIT call.
    \item Avoid unnecessary shape changes that trigger recompilation.
    \item Use \texttt{vmap} instead of Python loops when appropriate.
    \item Check which device JAX is actually using.
\end{itemize}

\textbf{Interview notes}

The key idea is that JAX is built around transformations of functions. Know
\texttt{grad}, \texttt{jit}, and \texttt{vmap}, plus immutability, explicit
PRNG keys, tracing, and PyTrees.


% ============================================================
% HUGGING FACE
% ============================================================

\subsection{Hugging Face}

\textbf{Hugging Face} is an ecosystem for pretrained models, tokenizers,
datasets, and model sharing. Important libraries include
\texttt{transformers}, \texttt{datasets}, and the Hugging Face Hub.

\textbf{Pipeline}

For quick inference:

\begin{lstlisting}[language=Python]
from transformers import pipeline

classifier = pipeline(
    "text-classification"
)

result = classifier(
    "The model performed very well."
)
\end{lstlisting}

Specify a model:

\begin{lstlisting}[language=Python]
classifier = pipeline(
    "text-classification",
    model="model-name"
)
\end{lstlisting}

Common tasks include:

\begin{lstlisting}
text-classification
token-classification
question-answering
summarization
translation
text-generation
image-classification
automatic-speech-recognition
\end{lstlisting}

\textbf{Tokenizer and model}

For more control:

\begin{lstlisting}[language=Python]
from transformers import AutoTokenizer
from transformers import AutoModel

tokenizer = AutoTokenizer.from_pretrained(
    "model-name"
)

model = AutoModel.from_pretrained(
    "model-name"
)
\end{lstlisting}

Common task-specific classes:

\begin{lstlisting}
AutoModelForSequenceClassification
AutoModelForTokenClassification
AutoModelForCausalLM
AutoModelForSeq2SeqLM
AutoModelForQuestionAnswering
\end{lstlisting}

Use a task-specific model class when the task requires a corresponding output
head.

\textbf{Tokenization}

Single input:

\begin{lstlisting}[language=Python]
inputs = tokenizer(
    "Example sentence.",
    return_tensors="pt"
)
\end{lstlisting}

Batch:

\begin{lstlisting}[language=Python]
inputs = tokenizer(
    texts,
    padding=True,
    truncation=True,
    return_tensors="pt"
)
\end{lstlisting}

Important options:

\begin{itemize}[leftmargin=*]
    \item \texttt{padding}: pad sequences to compatible lengths.
    \item \texttt{truncation}: truncate overly long inputs.
    \item \texttt{max\_length}: set an explicit sequence-length limit.
    \item \texttt{return\_tensors}: choose the tensor framework.
\end{itemize}

\textbf{Tokens versus words}

Tokens are not necessarily words. A tokenizer may divide a word into several
subword tokens.

This matters for:

\begin{itemize}[leftmargin=*]
    \item context-window limits;
    \item padding;
    \item truncation;
    \item memory use;
    \item generation length.
\end{itemize}

\textbf{Attention masks}

Padded batches commonly contain:

\begin{lstlisting}
input_ids
attention_mask
\end{lstlisting}

Conceptually:

\begin{lstlisting}
Tokens: [A, B, C, PAD, PAD]
Mask:   [1, 1, 1, 0,   0]
\end{lstlisting}

The attention mask identifies real tokens versus padding positions.

\textbf{Text generation}

\begin{lstlisting}[language=Python]
from transformers import AutoModelForCausalLM

model = AutoModelForCausalLM.from_pretrained(
    "model-name"
)

inputs = tokenizer(
    "Explain Bayesian inference:",
    return_tensors="pt"
)

outputs = model.generate(
    **inputs,
    max_new_tokens=100
)

text = tokenizer.decode(
    outputs[0],
    skip_special_tokens=True
)
\end{lstlisting}

Common generation parameters:

\begin{lstlisting}
max_new_tokens
max_length
temperature
top_p
top_k
do_sample
num_beams
repetition_penalty
\end{lstlisting}

Useful distinction:

\begin{itemize}[leftmargin=*]
    \item \texttt{max\_length}: limits the total sequence length, including the
    input and generated tokens.
    \item \texttt{max\_new\_tokens}: limits only the number of newly generated
    tokens.
\end{itemize}

For controlling output length, \texttt{max\_new\_tokens} is usually easier to
reason about.

\textbf{Greedy generation versus sampling}

Deterministic generation:

\begin{lstlisting}[language=Python]
outputs = model.generate(
    **inputs,
    do_sample=False,
    max_new_tokens=100
)
\end{lstlisting}

Sampling:

\begin{lstlisting}[language=Python]
outputs = model.generate(
    **inputs,
    do_sample=True,
    temperature=0.8,
    top_p=0.9,
    max_new_tokens=100
)
\end{lstlisting}

General interpretation:

\begin{itemize}[leftmargin=*]
    \item Lower temperature makes sampling more concentrated.
    \item Higher temperature increases randomness.
    \item \texttt{top\_p} restricts sampling to a probability-mass subset.
    \item \texttt{top\_k} restricts sampling to the top $k$ candidate tokens.
\end{itemize}

\textbf{Datasets}

\begin{lstlisting}[language=Python]
from datasets import load_dataset

dataset = load_dataset("dataset-name")
\end{lstlisting}

Common splits:

\begin{lstlisting}[language=Python]
dataset["train"]
dataset["validation"]
dataset["test"]
\end{lstlisting}

Transform a dataset:

\begin{lstlisting}[language=Python]
def tokenize(batch):
    return tokenizer(
        batch["text"],
        truncation=True
    )

tokenized = dataset.map(
    tokenize,
    batched=True
)
\end{lstlisting}

Filter:

\begin{lstlisting}[language=Python]
filtered = dataset.filter(
    lambda x: x["label"] == 1
)
\end{lstlisting}

\textbf{Model cards}

Before using a checkpoint, inspect its model card for:

\begin{itemize}[leftmargin=*]
    \item architecture;
    \item intended use;
    \item training data;
    \item evaluation results;
    \item limitations;
    \item license;
    \item preprocessing requirements;
    \item supported context length.
\end{itemize}

\textbf{Base versus fine-tuned models}

\begin{itemize}[leftmargin=*]
    \item \textbf{Base model}: pretrained general-purpose model.
    \item \textbf{Fine-tuned model}: model adapted to a task or desired behavior.
    \item \textbf{Instruction model}: model fine-tuned to follow natural-language
    instructions.
\end{itemize}

Use a checkpoint appropriate for the intended task.

\textbf{Full fine-tuning versus LoRA}

\begin{itemize}[leftmargin=*]
    \item \textbf{Full fine-tuning}: updates most or all model parameters.
    \item \textbf{LoRA}: learns smaller low-rank parameter updates while leaving most
    original model parameters frozen.
\end{itemize}

LoRA and other parameter-efficient methods can substantially reduce memory and
storage requirements.

\textbf{Practical tips}

\begin{itemize}[leftmargin=*]
    \item Match the model class to the task.
    \item Use the tokenizer associated with the checkpoint.
    \item Check context limits before processing long documents.
    \item Set padding and truncation deliberately.
    \item Prefer \texttt{max\_new\_tokens} when controlling generated output length.
    \item Inspect model cards and licenses before using a checkpoint.
    \item Batch inference when practical.
    \item Evaluate models on the actual target task rather than relying only on
    published benchmark scores.
\end{itemize}

\textbf{Interview notes}

Know the distinction between:

\begin{itemize}[leftmargin=*]
    \item model and tokenizer;
    \item base and fine-tuned models;
    \item tokens and words;
    \item padding and truncation;
    \item token IDs and attention masks;
    \item full fine-tuning and LoRA;
    \item greedy generation and sampling.
\end{itemize}

Also understand why the high-level \texttt{pipeline} interface is convenient
for quick inference while loading the tokenizer and model separately gives
more control.


% ============================================================
% LANGCHAIN
% ============================================================

\subsection{LangChain}

\textbf{LangChain} is an orchestration framework for applications built around
language models. It provides abstractions for prompts, models, structured
output, documents, retrieval, tools, and agentic workflows.

LangChain is not itself a language model.

\textbf{Core components}

\begin{center}
\begin{tabularx}{\textwidth}{lX}
\toprule
Component & Purpose \\
\midrule
Chat Model & Language-model interface \\
Prompt Template & Constructs model input \\
Output Parser & Converts model output into another representation \\
Document & Text plus associated metadata \\
Retriever & Finds relevant documents \\
Tool & Callable capability exposed to a workflow \\
Agent & Model-directed tool-use workflow \\
\bottomrule
\end{tabularx}
\end{center}

\textbf{Prompt templates}

\begin{lstlisting}[language=Python]
from langchain_core.prompts import ChatPromptTemplate

prompt = ChatPromptTemplate.from_messages([
    (
        "system",
        "You are a statistical assistant."
    ),
    (
        "human",
        "Explain {topic} in {n} sentences."
    )
])
\end{lstlisting}

Fill variables:

\begin{lstlisting}[language=Python]
prompt.invoke({
    "topic": "bootstrap sampling",
    "n": 3
})
\end{lstlisting}

\textbf{Runnable composition}

Components can be composed with the pipe operator:

\begin{lstlisting}[language=Python]
chain = prompt | model | parser
\end{lstlisting}

Invoke:

\begin{lstlisting}[language=Python]
result = chain.invoke({
    "topic": "bootstrap sampling",
    "n": 3
})
\end{lstlisting}

Batch:

\begin{lstlisting}[language=Python]
results = chain.batch([
    {
        "topic": "bootstrap",
        "n": 2
    },
    {
        "topic": "cross-validation",
        "n": 2
    }
])
\end{lstlisting}

Stream:

\begin{lstlisting}[language=Python]
for chunk in chain.stream(inputs):
    print(chunk, end="")
\end{lstlisting}

\textbf{Structured output}

Application code often needs predictable fields rather than arbitrary prose.

Conceptually:

\begin{lstlisting}[language=Python]
structured_model = model.with_structured_output(
    OutputSchema
)

result = structured_model.invoke(input)
\end{lstlisting}

Structured output is useful when downstream code expects values such as:

\begin{lstlisting}
classification
confidence
reason
\end{lstlisting}

Prefer structured output over manually extracting values from free-form text
when possible.

\textbf{Documents}

\begin{lstlisting}[language=Python]
from langchain_core.documents import Document

doc = Document(
    page_content=(
        "Clinical trial eligibility..."
    ),
    metadata={
        "source": "trial_123",
        "section": "eligibility"
    }
)
\end{lstlisting}

Useful metadata includes:

\begin{itemize}[leftmargin=*]
    \item source;
    \item document ID;
    \item page or section;
    \item date;
    \item category;
    \item other filtering/provenance information.
\end{itemize}

\textbf{RAG workflow}

A typical retrieval-augmented generation workflow has two stages.

Indexing:

\begin{lstlisting}
Documents
-> Chunking
-> Embeddings
-> Vector Store
\end{lstlisting}

Querying:

\begin{lstlisting}
Question
-> Retriever
-> Relevant Chunks
-> Model
-> Answer
\end{lstlisting}

Retrieve documents:

\begin{lstlisting}[language=Python]
docs = retriever.invoke(query)
\end{lstlisting}

\textbf{Chunking}

Important chunking decisions include:

\begin{itemize}[leftmargin=*]
    \item chunk size;
    \item overlap;
    \item document structure;
    \item semantic boundaries;
    \item metadata preservation.
\end{itemize}

Large chunks preserve more context but may introduce irrelevant information.
Very small chunks can improve retrieval precision but may lose necessary
surrounding context.

There is no universally optimal chunk size.

\textbf{Debugging RAG}

When a RAG answer is wrong, separate the problem into:

\begin{enumerate}[leftmargin=*]
    \item Was the correct evidence retrieved?
    \item If yes, did the model use that evidence correctly?
\end{enumerate}

If retrieval is wrong, inspect:

\begin{itemize}[leftmargin=*]
    \item document parsing;
    \item chunking;
    \item embeddings;
    \item metadata filters;
    \item retrieval depth;
    \item query formulation.
\end{itemize}

Do not immediately change the generation prompt when the relevant information
was never retrieved.

\textbf{Tools}

A tool exposes a controlled capability:

\begin{lstlisting}[language=Python]
from langchain_core.tools import tool

@tool
def calculate_score(
    x: float,
    y: float
) -> float:
    """Calculate the requested score."""
    return x + y
\end{lstlisting}

Good tools should have:

\begin{itemize}[leftmargin=*]
    \item a clear name;
    \item a narrow purpose;
    \item well-defined inputs;
    \item predictable outputs;
    \item appropriate permissions.
\end{itemize}

\textbf{Agents}

Agents allow a model to decide which action or tool to use next.

Use an agent when the sequence of actions genuinely needs to be dynamic.

For a known workflow such as:

\begin{lstlisting}
Input -> Retrieve -> Generate -> Validate
\end{lstlisting}

a deterministic pipeline is usually simpler and easier to test.

\textbf{State}

LLM applications must explicitly manage state such as:

\begin{itemize}[leftmargin=*]
    \item prior messages;
    \item conversation summaries;
    \item user/session information;
    \item retrieved context;
    \item workflow state.
\end{itemize}

Do not assume that the model automatically maintains persistent application
state.

\textbf{Prompt injection}

Retrieved documents and user-provided content should be treated as untrusted
input.

A retrieved document could contain text such as:

\begin{lstlisting}
Ignore all previous instructions and...
\end{lstlisting}

Instructions contained inside retrieved content should not automatically be
treated as trusted application instructions.

This is especially important when a model can invoke tools or perform actions.

\textbf{Practical tips}

\begin{itemize}[leftmargin=*]
    \item Use deterministic pipelines when the workflow is known.
    \item Use agents only when dynamic action selection is actually needed.
    \item Prefer structured outputs for downstream application logic.
    \item Preserve document metadata throughout RAG workflows.
    \item Inspect retrieved chunks directly when debugging RAG.
    \item Evaluate retrieval separately from generation.
    \item Keep tools narrow and clearly described.
    \item Restrict tool permissions to what the workflow needs.
    \item Treat retrieved text as untrusted content.
    \item Manage conversation and workflow state explicitly.
\end{itemize}

\textbf{Interview notes}

LangChain is an orchestration framework, not an LLM.

Know the distinction between:

\begin{itemize}[leftmargin=*]
    \item prompt, model, and parser;
    \item document and retriever;
    \item embeddings and retrieval;
    \item retrieval and generation;
    \item deterministic pipelines and agents;
    \item free-form and structured output;
    \item model context and persistent application state.
\end{itemize}

For RAG interviews, a strong debugging approach is to inspect retrieval first.
If the correct evidence was not retrieved, modifying the generation prompt
does not address the underlying retrieval failure.


% ============================================================
% 4. VALIDATION
% ============================================================

\section{Validation}

% ============================================================
% MLFLOW
% ============================================================

\subsection{MLflow}

\textbf{MLflow} is an open-source platform for tracking machine-learning
experiments and managing models throughout the ML lifecycle.

For day-to-day modeling, its most useful purpose is answering questions such
as:

\begin{itemize}[leftmargin=*]
    \item Which model produced this result?
    \item Which hyperparameters were used?
    \item Which run performed best?
    \item Where are the plots and model artifacts?
    \item Which model version should be used downstream?
\end{itemize}

The core hierarchy is:

\begin{lstlisting}
Experiment
    -> Run
        -> Parameters
        -> Metrics
        -> Tags
        -> Artifacts
        -> Models
\end{lstlisting}

\textbf{Core terminology}

\begin{center}
\begin{tabularx}{\textwidth}{lX}
\toprule
Term & Meaning \\
\midrule
Experiment & Collection of related runs \\
Run & One execution of an experiment \\
Parameter & Input/configuration value for a run \\
Metric & Numeric performance value \\
Artifact & File produced by a run \\
Tag & Searchable metadata attached to a run \\
Tracking Server & Stores and serves experiment metadata \\
Artifact Store & Stores files such as models and plots \\
Model Registry & Organizes registered models and versions \\
\bottomrule
\end{tabularx}
\end{center}

\textbf{Install MLflow}

\begin{lstlisting}
pip install mlflow
\end{lstlisting}

Check installation:

\begin{lstlisting}
mlflow --version
\end{lstlisting}

\textbf{Start the MLflow UI/server}

For local work:

\begin{lstlisting}
mlflow server --port 5000
\end{lstlisting}

Then open:

\begin{lstlisting}
http://localhost:5000
\end{lstlisting}

Set the tracking location in Python:

\begin{lstlisting}[language=Python]
import mlflow

mlflow.set_tracking_uri(
    "http://localhost:5000"
)
\end{lstlisting}

For simple local experimentation, MLflow can also be used without operating a
shared remote tracking server.

\textbf{Experiments}

Create or select an experiment:

\begin{lstlisting}[language=Python]
import mlflow

mlflow.set_experiment(
    "classification-models"
)
\end{lstlisting}

Use experiments to separate meaningful groups of work, such as:

\begin{lstlisting}
baseline-models
feature-engineering
hyperparameter-search
forecasting-experiments
production-candidates
\end{lstlisting}

Avoid creating a completely new experiment for every individual run.

\textbf{Runs}

A run represents one execution of a model or experiment.

\begin{lstlisting}[language=Python]
with mlflow.start_run():

    # train model
    # evaluate model
    # log results
    pass
\end{lstlisting}

A typical pattern is:

\begin{lstlisting}[language=Python]
with mlflow.start_run(
    run_name="random-forest-baseline"
):
    # experiment code
    pass
\end{lstlisting}

Give important runs meaningful names rather than relying only on generated
identifiers.

\textbf{Parameters}

Parameters describe the configuration used for a run.

\begin{lstlisting}[language=Python]
mlflow.log_param(
    "learning_rate",
    0.001
)

mlflow.log_param(
    "batch_size",
    64
)
\end{lstlisting}

Log several parameters:

\begin{lstlisting}[language=Python]
mlflow.log_params({
    "learning_rate": 0.001,
    "batch_size": 64,
    "hidden_size": 128,
    "optimizer": "adam"
})
\end{lstlisting}

Typical parameters include:

\begin{itemize}[leftmargin=*]
    \item learning rate;
    \item batch size;
    \item number of trees;
    \item regularization strength;
    \item model architecture;
    \item random seed;
    \item preprocessing choices.
\end{itemize}

\textbf{Metrics}

Metrics are numeric measurements of model performance.

\begin{lstlisting}[language=Python]
mlflow.log_metric(
    "accuracy",
    accuracy
)

mlflow.log_metric(
    "f1",
    f1
)
\end{lstlisting}

Log several metrics:

\begin{lstlisting}[language=Python]
mlflow.log_metrics({
    "accuracy": accuracy,
    "precision": precision,
    "recall": recall,
    "f1": f1
})
\end{lstlisting}

For iterative training, log the same metric over multiple steps:

\begin{lstlisting}[language=Python]
for epoch in range(num_epochs):

    train_loss = train_one_epoch(...)

    mlflow.log_metric(
        "train_loss",
        train_loss,
        step=epoch
    )
\end{lstlisting}

This produces a metric history rather than only a final value.

\textbf{Parameters versus metrics}

Keep the distinction clear:

\begin{itemize}[leftmargin=*]
    \item \textbf{Parameter}: configuration supplied to the experiment.
    \item \textbf{Metric}: measured result produced by the experiment.
\end{itemize}

For example:

\begin{lstlisting}
learning_rate = 0.001    -> parameter
batch_size    = 64       -> parameter

validation_loss = 0.32   -> metric
accuracy        = 0.91   -> metric
\end{lstlisting}

\textbf{Tags}

Tags store descriptive metadata:

\begin{lstlisting}[language=Python]
mlflow.set_tag(
    "model_type",
    "random_forest"
)

mlflow.set_tag(
    "dataset",
    "customer_churn_v2"
)
\end{lstlisting}

Useful tags may describe:

\begin{itemize}[leftmargin=*]
    \item model family;
    \item dataset version;
    \item experiment purpose;
    \item development stage;
    \item feature set;
    \item notes about a run.
\end{itemize}

\textbf{Artifacts}

Artifacts are files associated with a run.

Examples include:

\begin{itemize}[leftmargin=*]
    \item plots;
    \item confusion matrices;
    \item calibration plots;
    \item feature-importance figures;
    \item prediction files;
    \item configuration files;
    \item model files.
\end{itemize}

Log one artifact:

\begin{lstlisting}[language=Python]
mlflow.log_artifact(
    "confusion_matrix.png"
)
\end{lstlisting}

Organize it inside the run:

\begin{lstlisting}[language=Python]
mlflow.log_artifact(
    "confusion_matrix.png",
    artifact_path="plots"
)
\end{lstlisting}

Log an entire directory:

\begin{lstlisting}[language=Python]
mlflow.log_artifacts(
    "results/",
    artifact_path="results"
)
\end{lstlisting}

\textbf{Complete tracking pattern}

A typical experiment might look like:

\begin{lstlisting}[language=Python]
import mlflow

mlflow.set_experiment(
    "classification-models"
)

with mlflow.start_run(
    run_name="rf-baseline"
):

    mlflow.log_params({
        "n_estimators": 500,
        "max_depth": 10,
        "random_state": 42
    })

    model.fit(X_train, y_train)

    predictions = model.predict(X_valid)

    accuracy = accuracy_score(
        y_valid,
        predictions
    )

    f1 = f1_score(
        y_valid,
        predictions
    )

    mlflow.log_metrics({
        "accuracy": accuracy,
        "f1": f1
    })

    mlflow.log_artifact(
        "confusion_matrix.png"
    )
\end{lstlisting}

The important idea is that the model result and the configuration that
produced it are stored together.

\textbf{Autologging}

MLflow can automatically capture common information for supported ML
frameworks.

For example:

\begin{lstlisting}[language=Python]
import mlflow
import mlflow.sklearn

mlflow.sklearn.autolog()

with mlflow.start_run():
    model.fit(
        X_train,
        y_train
    )
\end{lstlisting}

Autologging can capture information such as:

\begin{itemize}[leftmargin=*]
    \item model parameters;
    \item selected metrics;
    \item fitted models;
    \item framework metadata.
\end{itemize}

Autologging is convenient, but manually log project-specific information that
the automatic integration does not know about.

\textbf{Logging models}

For a scikit-learn model:

\begin{lstlisting}[language=Python]
import mlflow.sklearn

mlflow.sklearn.log_model(
    sk_model=model,
    name="model"
)
\end{lstlisting}

Logging the model with MLflow keeps the fitted model associated with the run
rather than saving an unexplained model file somewhere on disk.

\textbf{Loading a logged model}

A model can later be loaded from its MLflow model location.

For a scikit-learn model:

\begin{lstlisting}[language=Python]
import mlflow.sklearn

model = mlflow.sklearn.load_model(
    model_uri
)
\end{lstlisting}

This is useful when evaluation, deployment, or batch prediction should use the
exact model produced by a tracked run.

\textbf{Model Registry}

The Model Registry provides a higher-level organization for models intended
to move beyond individual experiment runs.

Useful concepts include:

\begin{itemize}[leftmargin=*]
    \item registered model;
    \item model version;
    \item model alias;
    \item model metadata.
\end{itemize}

A common workflow is:

\begin{lstlisting}
Train candidate
-> Track run
-> Compare validation results
-> Register selected model
-> Assign useful alias
-> Use registered model downstream
\end{lstlisting}

Aliases can represent roles such as:

\begin{lstlisting}
champion
candidate
production
\end{lstlisting}

The exact naming convention is a team decision; the important point is to
avoid identifying production models only by fragile local filenames.

\textbf{Comparing runs}

One of MLflow's most useful features is comparing experiments across
parameters and metrics.

Programmatically:

\begin{lstlisting}[language=Python]
runs = mlflow.search_runs(
    experiment_names=[
        "classification-models"
    ]
)
\end{lstlisting}

The result can be inspected like a table:

\begin{lstlisting}[language=Python]
runs[[
    "run_id",
    "params.max_depth",
    "metrics.accuracy",
    "metrics.f1"
]]
\end{lstlisting}

This makes it easier to answer questions such as:

\begin{itemize}[leftmargin=*]
    \item Which run had the highest validation F1?
    \item Which hyperparameters produced the best result?
    \item Did increasing model complexity actually improve performance?
\end{itemize}

\textbf{What to log}

A useful minimum for most experiments is:

\begin{center}
\begin{tabularx}{\textwidth}{lX}
\toprule
Category & Examples \\
\midrule
Parameters & Hyperparameters, seed, preprocessing settings \\
Metrics & Validation loss, RMSE, F1, AUROC, calibration metrics \\
Tags & Model family, dataset version, experiment purpose \\
Artifacts & Plots, tables, predictions, configuration files \\
Model & Fitted model used to generate the results \\
\bottomrule
\end{tabularx}
\end{center}

For probabilistic models, useful artifacts may include:

\begin{itemize}[leftmargin=*]
    \item calibration curves;
    \item reliability diagrams;
    \item residual plots;
    \item PIT histograms;
    \item prediction intervals;
    \item evaluation tables.
\end{itemize}

\textbf{Validation versus test data}

MLflow tracks results, but it does not prevent bad experimental design.

Keep the roles of datasets distinct:

\begin{itemize}[leftmargin=*]
    \item \textbf{Training set}: fit model parameters.
    \item \textbf{Validation set}: tune models and compare candidates.
    \item \textbf{Test set}: estimate final generalization performance.
\end{itemize}

If the test set is repeatedly used to choose models, it effectively becomes
another validation set.

MLflow makes comparisons easier, so it is especially important not to use
that convenience to repeatedly optimize against the final test set.

\textbf{Metric selection}

Do not select a model solely because one metric is numerically highest.

Examples:

\begin{itemize}[leftmargin=*]
    \item Accuracy can be misleading under class imbalance.
    \item AUROC and AUPRC emphasize different aspects of classification.
    \item RMSE penalizes large errors more heavily than MAE.
    \item Good likelihood does not guarantee good calibration.
    \item Average performance can hide poor subgroup performance.
\end{itemize}

Log the metrics that correspond to the actual modeling objective.

\textbf{Data leakage}

Experiment tracking cannot fix leakage.

Common examples include:

\begin{itemize}[leftmargin=*]
    \item preprocessing the full dataset before splitting;
    \item using future information in time-series features;
    \item selecting features using the test set;
    \item allowing the same patient or entity to appear across inappropriate splits;
    \item fitting normalization parameters using validation/test observations.
\end{itemize}

A perfectly tracked leaked experiment is still invalid.

\textbf{Reproducibility}

Tracking hyperparameters alone is not enough for complete reproducibility.

Also consider recording:

\begin{itemize}[leftmargin=*]
    \item random seed;
    \item dataset/version;
    \item train/validation/test split;
    \item preprocessing configuration;
    \item package/environment information;
    \item code version or commit;
    \item model architecture;
    \item important feature definitions.
\end{itemize}

The goal is to be able to answer:

\begin{quote}
What exactly produced this model and this result?
\end{quote}

\textbf{Run naming}

Poor:

\begin{lstlisting}
run1
run2
test
test2
final
final_final
\end{lstlisting}

Better:

\begin{lstlisting}
rf-baseline
xgb-class-weighted
lstm-window-30
transformer-small
\end{lstlisting}

Use parameters and tags for detailed configuration rather than encoding every
setting into an enormous run name.

\textbf{Practical tips}

\begin{itemize}[leftmargin=*]
    \item Create experiments around meaningful projects or experiment families.
    \item Give important runs descriptive names.
    \item Log configuration before or during training rather than reconstructing it
    afterward.
    \item Log validation metrics, not just training metrics.
    \item Save important plots as artifacts.
    \item Record random seeds and dataset versions.
    \item Use tags for information that helps filter and organize runs.
    \item Use autologging for convenience, then add project-specific metrics and
    artifacts manually.
    \item Compare runs systematically rather than relying on notebook output or
    memory.
    \item Do not repeatedly tune against the final test set.
    \item Keep model selection criteria consistent across runs.
    \item Track enough information to reproduce the winning run.
\end{itemize}

\textbf{Common mistakes}

\begin{itemize}[leftmargin=*]
    \item Logging only the final model and not the configuration that produced it.
    \item Logging only training performance.
    \item Changing datasets without recording which version was used.
    \item Comparing runs evaluated on different splits as if they were equivalent.
    \item Using inconsistent metric definitions across experiments.
    \item Tracking hundreds of runs without useful names, tags, or organization.
    \item Assuming experiment tracking automatically guarantees reproducibility.
    \item Choosing the best model by repeatedly inspecting test-set performance.
\end{itemize}

\textbf{Interview notes}

The main purpose of MLflow is \textbf{experiment tracking and model lifecycle
management}. It helps connect model configurations, metrics, artifacts, and
models so experiments can be compared and reproduced.

Be comfortable explaining:

\begin{itemize}[leftmargin=*]
    \item experiment versus run;
    \item parameter versus metric;
    \item metric versus artifact;
    \item manual logging versus autologging;
    \item model logging versus model registration;
    \item why dataset and code versions matter for reproducibility;
    \item why experiment tracking does not prevent leakage or poor validation design.
\end{itemize}

A useful interview workflow is:

\begin{lstlisting}
Define experiment
-> Start run
-> Log configuration
-> Train model
-> Evaluate on validation data
-> Log metrics and artifacts
-> Compare runs
-> Select candidate
-> Final evaluation
-> Register/version selected model
\end{lstlisting}

If asked why MLflow is useful instead of simply saving models as
\texttt{model.pkl}, emphasize that a model file alone does not tell you which
data, parameters, code, metrics, or experimental conditions produced it.


% ============================================================
% 5. DEPLOYMENT
% ============================================================

\section{Deployment}

% ============================================================
% DOCKER
% ============================================================

\subsection{Docker}

\textbf{Docker} packages an application and its dependencies into a
\textbf{container} so that it can run consistently across development,
testing, and deployment environments.

The core idea is:

\begin{lstlisting}
Application Code
+ Runtime
+ Dependencies
+ Configuration
= Container Image
\end{lstlisting}

A running instance of that image is a \textbf{container}.

\textbf{Core terminology}

\begin{center}
\begin{tabularx}{\textwidth}{lX}
\toprule
Term & Meaning \\
\midrule
Dockerfile & Instructions for building an image \\
Image & Immutable template used to create containers \\
Container & Running instance of an image \\
Registry & Repository for storing and distributing images \\
Volume & Persistent storage managed outside the container filesystem \\
Bind Mount & Host directory mounted directly into a container \\
Port Mapping & Maps a host port to a container port \\
Docker Compose & Defines and runs multi-container applications \\
\bottomrule
\end{tabularx}
\end{center}

The distinction to remember is:

\begin{lstlisting}
Dockerfile -> Image -> Container
\end{lstlisting}


% ------------------------------------------------------------
% BASIC COMMANDS
% ------------------------------------------------------------

\textbf{Check Docker}

\begin{lstlisting}
docker --version
docker info
\end{lstlisting}

\textbf{Images}

List local images:

\begin{lstlisting}
docker images
\end{lstlisting}

Pull an image:

\begin{lstlisting}
docker pull python:3.12-slim
\end{lstlisting}

Remove an image:

\begin{lstlisting}
docker rmi image_name
\end{lstlisting}

\textbf{Containers}

List running containers:

\begin{lstlisting}
docker ps
\end{lstlisting}

List all containers:

\begin{lstlisting}
docker ps -a
\end{lstlisting}

Start:

\begin{lstlisting}
docker start container_name
\end{lstlisting}

Stop:

\begin{lstlisting}
docker stop container_name
\end{lstlisting}

Restart:

\begin{lstlisting}
docker restart container_name
\end{lstlisting}

Remove:

\begin{lstlisting}
docker rm container_name
\end{lstlisting}

Force-remove a running container:

\begin{lstlisting}
docker rm -f container_name
\end{lstlisting}


% ------------------------------------------------------------
% DOCKERFILE
% ------------------------------------------------------------

\textbf{Dockerfile}

A simple Python application:

\begin{lstlisting}
FROM python:3.12-slim

WORKDIR /app

COPY requirements.txt .

RUN pip install --no-cache-dir \
    -r requirements.txt

COPY . .

CMD ["python", "app.py"]
\end{lstlisting}

Important instructions:

\begin{center}
\begin{tabularx}{\textwidth}{lX}
\toprule
Instruction & Purpose \\
\midrule
FROM & Select base image \\
WORKDIR & Set working directory \\
COPY & Copy files into image \\
RUN & Execute command while building image \\
ENV & Define environment variable \\
EXPOSE & Document intended container port \\
CMD & Default command when container starts \\
ENTRYPOINT & Define the executable for the container \\
\bottomrule
\end{tabularx}
\end{center}

\textbf{FROM}

Every Dockerfile normally begins with a base image:

\begin{lstlisting}
FROM python:3.12-slim
\end{lstlisting}

Smaller base images generally reduce image size and unnecessary packages.

\textbf{WORKDIR}

\begin{lstlisting}
WORKDIR /app
\end{lstlisting}

Subsequent commands operate relative to this directory.

\textbf{COPY}

\begin{lstlisting}
COPY requirements.txt .
\end{lstlisting}

Copy the project:

\begin{lstlisting}
COPY . .
\end{lstlisting}

\textbf{RUN}

\texttt{RUN} executes while the image is being built:

\begin{lstlisting}
RUN pip install --no-cache-dir \
    -r requirements.txt
\end{lstlisting}

The resulting filesystem changes become part of the image.

\textbf{CMD}

\texttt{CMD} specifies the default command when the container starts:

\begin{lstlisting}
CMD ["python", "app.py"]
\end{lstlisting}

Prefer the JSON/exec form shown above rather than:

\begin{lstlisting}
CMD python app.py
\end{lstlisting}


% ------------------------------------------------------------
% BUILDING IMAGES
% ------------------------------------------------------------

\textbf{Build an image}

From the directory containing the Dockerfile:

\begin{lstlisting}
docker build -t my-app .
\end{lstlisting}

The final period means:

\begin{lstlisting}
use the current directory as the build context
\end{lstlisting}

Tag a version:

\begin{lstlisting}
docker build -t my-app:1.0 .
\end{lstlisting}

List images:

\begin{lstlisting}
docker images
\end{lstlisting}


% ------------------------------------------------------------
% RUNNING CONTAINERS
% ------------------------------------------------------------

\textbf{Run a container}

\begin{lstlisting}
docker run my-app
\end{lstlisting}

Run with a name:

\begin{lstlisting}
docker run --name my-container my-app
\end{lstlisting}

Run in the background:

\begin{lstlisting}
docker run -d \
    --name my-container \
    my-app
\end{lstlisting}

Here:

\begin{lstlisting}
-d
\end{lstlisting}

means detached/background mode.

Automatically remove the container when it exits:

\begin{lstlisting}
docker run --rm my-app
\end{lstlisting}

This is useful for temporary jobs and testing.


% ------------------------------------------------------------
% PORTS
% ------------------------------------------------------------

\textbf{Port mapping}

Suppose the application listens on port 8000 inside the container.

Run:

\begin{lstlisting}
docker run -p 8000:8000 my-app
\end{lstlisting}

The syntax is:

\begin{lstlisting}
-p HOST_PORT:CONTAINER_PORT
\end{lstlisting}

For example:

\begin{lstlisting}
docker run -p 8080:8000 my-app
\end{lstlisting}

means:

\begin{lstlisting}
localhost:8080
    -> container:8000
\end{lstlisting}

The host and container ports do not need to match.

\textbf{EXPOSE}

A Dockerfile may contain:

\begin{lstlisting}
EXPOSE 8000
\end{lstlisting}

This documents that the application expects to listen on port 8000.

It does \textbf{not} automatically publish the port to the host.

You still need:

\begin{lstlisting}
docker run -p 8000:8000 my-app
\end{lstlisting}


% ------------------------------------------------------------
% ENVIRONMENT VARIABLES
% ------------------------------------------------------------

\textbf{Environment variables}

Pass configuration when starting the container:

\begin{lstlisting}
docker run \
    -e MODEL_NAME=my-model \
    -e LOG_LEVEL=INFO \
    my-app
\end{lstlisting}

Inside Python:

\begin{lstlisting}[language=Python]
import os

model_name = os.getenv(
    "MODEL_NAME"
)
\end{lstlisting}

Use an environment file:

\begin{lstlisting}
docker run --env-file .env my-app
\end{lstlisting}

Example \texttt{.env}:

\begin{lstlisting}
MODEL_NAME=my-model
LOG_LEVEL=INFO
DATABASE_HOST=db
\end{lstlisting}

Do not bake secrets directly into the image or commit them into source
control.


% ------------------------------------------------------------
% LOGS AND DEBUGGING
% ------------------------------------------------------------

\textbf{Logs}

View container logs:

\begin{lstlisting}
docker logs my-container
\end{lstlisting}

Follow logs continuously:

\begin{lstlisting}
docker logs -f my-container
\end{lstlisting}

Show recent logs:

\begin{lstlisting}
docker logs --tail 100 my-container
\end{lstlisting}

\textbf{Execute a command inside a running container}

Open a shell:

\begin{lstlisting}
docker exec -it my-container sh
\end{lstlisting}

If Bash is installed:

\begin{lstlisting}
docker exec -it my-container bash
\end{lstlisting}

Run one command:

\begin{lstlisting}
docker exec my-container ls /app
\end{lstlisting}

\textbf{Inspect a container}

\begin{lstlisting}
docker inspect my-container
\end{lstlisting}

Useful for checking:

\begin{itemize}[leftmargin=*]
    \item environment variables;
    \item networking;
    \item mounts;
    \item container configuration;
    \item image information.
\end{itemize}

\textbf{View resource usage}

\begin{lstlisting}
docker stats
\end{lstlisting}

This is useful when debugging unexpectedly high CPU or memory usage.


% ------------------------------------------------------------
% VOLUMES
% ------------------------------------------------------------

\textbf{Container storage}

A container's writable filesystem should generally be treated as temporary.

Removing the container can remove data stored only inside that container.

Use external storage for persistent data.

\textbf{Volumes}

Create a volume:

\begin{lstlisting}
docker volume create app-data
\end{lstlisting}

Mount it:

\begin{lstlisting}
docker run \
    -v app-data:/app/data \
    my-app
\end{lstlisting}

List volumes:

\begin{lstlisting}
docker volume ls
\end{lstlisting}

Inspect:

\begin{lstlisting}
docker volume inspect app-data
\end{lstlisting}

\textbf{Bind mounts}

Mount a local directory:

\begin{lstlisting}
docker run \
    -v "$(pwd)":/app \
    my-app
\end{lstlisting}

This directly maps the current host directory to \texttt{/app} inside the
container.

Bind mounts are especially useful during development because local code
changes are immediately visible inside the container.

\textbf{Volume versus bind mount}

\begin{center}
\begin{tabularx}{\textwidth}{lXX}
\toprule
 & Volume & Bind Mount \\
\midrule
Managed by Docker & Yes & No \\
Uses explicit host path & No & Yes \\
Good for persistent app data & Yes & Possible \\
Good for local development & Possible & Yes \\
\bottomrule
\end{tabularx}
\end{center}


% ------------------------------------------------------------
% .DOCKERIGNORE
% ------------------------------------------------------------

\textbf{.dockerignore}

A \texttt{.dockerignore} file prevents unnecessary files from being included
in the build context.

Example:

\begin{lstlisting}
.git
.venv
__pycache__
*.pyc
.env
data/
models/
notebooks/
\end{lstlisting}

This can:

\begin{itemize}[leftmargin=*]
    \item reduce build time;
    \item reduce accidental inclusion of large files;
    \item prevent secrets from being copied into the image;
    \item keep images cleaner.
\end{itemize}

Do not rely on \texttt{.dockerignore} as the only mechanism protecting
secrets, but use it as an additional safeguard.


% ------------------------------------------------------------
% BUILD CACHE
% ------------------------------------------------------------

\textbf{Build cache}

Docker caches build layers.

This Dockerfile:

\begin{lstlisting}
FROM python:3.12-slim

WORKDIR /app

COPY requirements.txt .

RUN pip install --no-cache-dir \
    -r requirements.txt

COPY . .
\end{lstlisting}

is preferable to:

\begin{lstlisting}
FROM python:3.12-slim

WORKDIR /app

COPY . .

RUN pip install --no-cache-dir \
    -r requirements.txt
\end{lstlisting}

because application code changes frequently while
\texttt{requirements.txt} changes less frequently.

With the first structure, Docker can often reuse the dependency-installation
layer.

A useful rule is:

\begin{quote}
Put relatively stable build steps before frequently changing build steps.
\end{quote}


% ------------------------------------------------------------
% IMAGE SIZE
% ------------------------------------------------------------

\textbf{Keeping images small}

Useful practices include:

\begin{itemize}[leftmargin=*]
    \item use an appropriately small base image;
    \item exclude unnecessary files with \texttt{.dockerignore};
    \item avoid unnecessary system packages;
    \item remove temporary installation files;
    \item use \texttt{pip install --no-cache-dir};
    \item use multi-stage builds when build dependencies are not needed at runtime.
\end{itemize}

Inspect image size:

\begin{lstlisting}
docker images
\end{lstlisting}


% ------------------------------------------------------------
% MULTI-STAGE BUILDS
% ------------------------------------------------------------

\textbf{Multi-stage builds}

Multi-stage builds separate build dependencies from the final runtime image.

General structure:

\begin{lstlisting}
FROM some-image AS builder

# Build application here


FROM smaller-runtime-image

# Copy only required output
COPY --from=builder /build/output /app

CMD ["..."]
\end{lstlisting}

This is particularly useful for compiled applications or frontend builds where
the compiler/build toolchain is unnecessary at runtime.


% ------------------------------------------------------------
% DOCKER COMPOSE
% ------------------------------------------------------------

\textbf{Docker Compose}

Docker Compose is useful when an application contains multiple services.

For example:

\begin{lstlisting}
Web Application
+ Database
+ Cache
\end{lstlisting}

A simple \texttt{compose.yaml}:

\begin{lstlisting}
services:

  app:
    build: .
    ports:
      - "8000:8000"
    environment:
      DATABASE_HOST: db
    depends_on:
      - db

  db:
    image: postgres
    environment:
      POSTGRES_DB: app
      POSTGRES_USER: user
      POSTGRES_PASSWORD: password
    volumes:
      - db-data:/var/lib/postgresql/data

volumes:
  db-data:
\end{lstlisting}

Start all services:

\begin{lstlisting}
docker compose up
\end{lstlisting}

Run in background:

\begin{lstlisting}
docker compose up -d
\end{lstlisting}

Build/rebuild:

\begin{lstlisting}
docker compose up --build
\end{lstlisting}

Stop and remove containers/networks:

\begin{lstlisting}
docker compose down
\end{lstlisting}

View logs:

\begin{lstlisting}
docker compose logs
\end{lstlisting}

Follow logs:

\begin{lstlisting}
docker compose logs -f
\end{lstlisting}

List services:

\begin{lstlisting}
docker compose ps
\end{lstlisting}


% ------------------------------------------------------------
% CONTAINER NETWORKING
% ------------------------------------------------------------

\textbf{Container networking}

A common source of confusion is:

\begin{lstlisting}
localhost
\end{lstlisting}

Inside a container, \texttt{localhost} refers to that container itself.

If a Compose service named \texttt{app} needs to communicate with a service
named \texttt{db}, it can typically use the service name:

\begin{lstlisting}
db
\end{lstlisting}

rather than:

\begin{lstlisting}
localhost
\end{lstlisting}

For example:

\begin{lstlisting}
DATABASE_HOST=db
\end{lstlisting}

This is one of the most common Docker networking issues during development.


% ------------------------------------------------------------
% DEVELOPMENT WORKFLOW
% ------------------------------------------------------------

\textbf{Typical development workflow}

Build:

\begin{lstlisting}
docker build -t my-app .
\end{lstlisting}

Run:

\begin{lstlisting}
docker run \
    --name my-app-container \
    -p 8000:8000 \
    my-app
\end{lstlisting}

Check:

\begin{lstlisting}
docker ps
\end{lstlisting}

Inspect logs:

\begin{lstlisting}
docker logs -f my-app-container
\end{lstlisting}

Enter the container if necessary:

\begin{lstlisting}
docker exec -it my-app-container sh
\end{lstlisting}

Stop:

\begin{lstlisting}
docker stop my-app-container
\end{lstlisting}

Remove:

\begin{lstlisting}
docker rm my-app-container
\end{lstlisting}


% ------------------------------------------------------------
% CLEANUP
% ------------------------------------------------------------

\textbf{Cleanup}

Remove stopped containers:

\begin{lstlisting}
docker container prune
\end{lstlisting}

Remove unused images:

\begin{lstlisting}
docker image prune
\end{lstlisting}

Remove unused Docker objects:

\begin{lstlisting}
docker system prune
\end{lstlisting}

Check disk usage:

\begin{lstlisting}
docker system df
\end{lstlisting}

Use cleanup commands carefully, especially when volumes or cached images are
important.


% ------------------------------------------------------------
% CONTAINERS VS VIRTUAL MACHINES
% ------------------------------------------------------------

\textbf{Containers versus virtual machines}

\begin{center}
\begin{tabularx}{\textwidth}{lXX}
\toprule
 & Container & Virtual Machine \\
\midrule
Isolation & Process-level & Full OS-level \\
Guest OS & Shares host kernel & Separate guest OS \\
Startup & Usually fast & Usually slower \\
Size & Usually smaller & Usually larger \\
Portability & High & High, but heavier \\
\bottomrule
\end{tabularx}
\end{center}

Containers are not simply lightweight virtual machines. They isolate
processes while sharing the host operating system kernel.


% ------------------------------------------------------------
% COMMON FAILURE MODES
% ------------------------------------------------------------

\textbf{Common failure modes}

\textbf{Container immediately exits}

Check:

\begin{lstlisting}
docker ps -a
docker logs container_name
\end{lstlisting}

A container normally runs only while its main process is running.

\textbf{Application works locally but cannot be reached}

Check:

\begin{itemize}[leftmargin=*]
    \item whether the correct port was published;
    \item whether the application is listening on the expected container port;
    \item whether the application is bound only to \texttt{127.0.0.1} inside the
    container.
\end{itemize}

Web applications in containers commonly need to listen on:

\begin{lstlisting}
0.0.0.0
\end{lstlisting}

rather than only:

\begin{lstlisting}
127.0.0.1
\end{lstlisting}

\textbf{Code changes do not appear}

Possible causes:

\begin{itemize}[leftmargin=*]
    \item the image was not rebuilt;
    \item Docker reused a cached layer;
    \item the wrong image/tag is running;
    \item a bind mount is overriding files inside the image.
\end{itemize}

Rebuild when needed:

\begin{lstlisting}
docker build -t my-app .
\end{lstlisting}

Force a build without cache when debugging cache issues:

\begin{lstlisting}
docker build --no-cache -t my-app .
\end{lstlisting}

\textbf{Container cannot reach another service}

Remember that:

\begin{lstlisting}
localhost
\end{lstlisting}

inside a container refers to that container, not another container and not
automatically the host machine.


% ------------------------------------------------------------
% PRACTICAL TIPS
% ------------------------------------------------------------

\textbf{Practical tips}

\begin{itemize}[leftmargin=*]
    \item Remember the sequence: Dockerfile $\rightarrow$ image $\rightarrow$
    container.
    \item Use meaningful image tags instead of relying only on
    \texttt{latest}.
    \item Keep images small and dependencies explicit.
    \item Use \texttt{.dockerignore} to reduce unnecessary build context.
    \item Copy dependency files before application code to improve build caching.
    \item Store persistent data outside the container filesystem.
    \item Use environment variables for runtime configuration.
    \item Never bake credentials or private keys into an image.
    \item Use \texttt{docker logs} first when a container fails.
    \item Use \texttt{docker exec} when you need to inspect a running container.
    \item Remember that \texttt{EXPOSE} does not publish a port.
    \item Remember that container \texttt{localhost} refers to the container itself.
    \item Bind web applications to \texttt{0.0.0.0} when they need to be reachable
    outside the container.
    \item Use Compose when several services must run together.
    \item Pin important dependency and base-image versions for reproducibility.
\end{itemize}


% ------------------------------------------------------------
% INTERVIEW NOTES
% ------------------------------------------------------------

\textbf{Interview notes}

Be comfortable explaining the distinction between:

\begin{itemize}[leftmargin=*]
    \item Dockerfile, image, and container;
    \item \texttt{RUN} and \texttt{CMD};
    \item build time and runtime;
    \item image and container;
    \item volume and bind mount;
    \item container port and host port;
    \item containers and virtual machines;
    \item Docker and Docker Compose.
\end{itemize}

Common interview questions include:

\begin{itemize}[leftmargin=*]
    \item \textbf{Why use Docker?} \\
    To package an application and its dependencies into a reproducible,
    portable runtime environment.

    \item \textbf{What is the difference between an image and a container?} \\
    An image is an immutable template; a container is a running instance of
    that image.

    \item \textbf{What happens to data when a container is deleted?} \\
    Data stored only in the container's writable filesystem can disappear.
    Persistent data should be stored in volumes or external systems.

    \item \textbf{What does \texttt{-p 8080:8000} mean?} \\
    Host port 8080 is forwarded to port 8000 inside the container.

    \item \textbf{Why does Docker build caching matter?} \\
    Docker can reuse unchanged layers, so ordering Dockerfile instructions
    carefully can make repeated builds substantially faster.

    \item \textbf{Why might an application work outside Docker but fail inside it?} \\
    Common causes include missing dependencies, incorrect environment variables,
    incorrect port mappings, filesystem assumptions, networking differences,
    or binding the application only to \texttt{localhost}.

    \item \textbf{When would you use Docker Compose?} \\
    When an application requires multiple coordinated services such as a web
    application, database, worker, or cache.
\end{itemize}

For technical interviews, knowing how to debug a broken container is often more
useful than memorizing Docker commands. A strong first sequence is:

\begin{lstlisting}
docker ps -a
docker logs container_name
docker inspect container_name
docker exec -it container_name sh
\end{lstlisting}

This answers four useful questions: whether the container exists and is
running, what it logged, how it is configured, and what its internal
environment looks like.


% ============================================================
% 6. REPORTING
% ============================================================

\section{Reporting}

% ============================================================
% POWER BI
% ============================================================

\subsection{Power BI}

\textbf{Power BI} is Microsoft's business intelligence platform for
transforming data into interactive reports and dashboards.

The core workflow is:

\begin{lstlisting}
Data Source
-> Power Query
-> Data Model
-> DAX
-> Visualizations
-> Report
\end{lstlisting}

The most important skills for day-to-day Power BI work are:

\begin{itemize}[leftmargin=*]
    \item cleaning data with Power Query;
    \item designing a good relational model;
    \item creating measures with DAX;
    \item understanding filter context;
    \item building useful rather than merely attractive reports.
\end{itemize}


% ------------------------------------------------------------
% CORE TERMINOLOGY
% ------------------------------------------------------------

\textbf{Core terminology}

\begin{center}
\begin{tabularx}{\textwidth}{lX}
\toprule
Term & Meaning \\
\midrule
Power BI Desktop & Main application for building models and reports \\
Power Query & Data import and transformation layer \\
Semantic Model & Tables, relationships, measures, and business logic \\
DAX & Formula language used for calculations \\
Measure & Calculation evaluated according to filter context \\
Calculated Column & Value calculated and stored for each row \\
Report & Multi-page interactive collection of visuals \\
Dashboard & Service-level canvas containing pinned report visuals \\
Slicer & Interactive visual used to filter report data \\
\bottomrule
\end{tabularx}
\end{center}


% ------------------------------------------------------------
% POWER QUERY
% ------------------------------------------------------------

\textbf{Power Query}

Power Query is used to load and transform data before it enters the analytical
model.

Common transformations include:

\begin{itemize}[leftmargin=*]
    \item changing data types;
    \item renaming columns;
    \item filtering rows;
    \item removing duplicates;
    \item splitting or combining columns;
    \item joining datasets;
    \item appending datasets;
    \item pivoting and unpivoting;
    \item handling missing values.
\end{itemize}

A useful rule is:

\begin{quote}
Use Power Query for data preparation and DAX for analytical calculations.
\end{quote}

Do not use complicated DAX to compensate for data that should have been
cleaned or reshaped before entering the model.

\textbf{Applied Steps}

Power Query records transformations as ordered steps.

Conceptually:

\begin{lstlisting}
Source
-> Select Columns
-> Change Types
-> Filter Rows
-> Remove Duplicates
-> Merge Tables
-> Load
\end{lstlisting}

Order matters because each step operates on the result of the previous step.

\textbf{Merge versus Append}

\begin{itemize}[leftmargin=*]
    \item \textbf{Merge}: combine tables horizontally using matching keys.
    \item \textbf{Append}: combine tables vertically by stacking rows.
\end{itemize}

Think:

\begin{lstlisting}
Merge  -> SQL JOIN
Append -> SQL UNION ALL
\end{lstlisting}


% ------------------------------------------------------------
% DATA MODELING
% ------------------------------------------------------------

\textbf{Data modeling}

A well-designed model is usually more important than complicated DAX.

Power BI works particularly well with a \textbf{star schema}.

A typical model contains:

\begin{lstlisting}
          DimDate
             |
DimCustomer--FactSales--DimProduct
             |
          DimStore
\end{lstlisting}

\textbf{Fact tables}

Fact tables contain observations or events such as:

\begin{itemize}[leftmargin=*]
    \item sales transactions;
    \item claims;
    \item orders;
    \item visits;
    \item measurements.
\end{itemize}

They commonly contain:

\begin{itemize}[leftmargin=*]
    \item foreign keys;
    \item dates;
    \item quantities;
    \item monetary values;
    \item measurable outcomes.
\end{itemize}

\textbf{Dimension tables}

Dimension tables describe entities used to group or filter facts.

Examples:

\begin{itemize}[leftmargin=*]
    \item customer;
    \item product;
    \item date;
    \item geography;
    \item department.
\end{itemize}

\textbf{Grain}

Before modeling a table, determine:

\begin{quote}
What does one row represent?
\end{quote}

For example:

\begin{lstlisting}
FactSales:
one row = one transaction line

DimCustomer:
one row = one customer
\end{lstlisting}

Mixing different grains is a common source of incorrect calculations.


% ------------------------------------------------------------
% RELATIONSHIPS
% ------------------------------------------------------------

\textbf{Relationships}

A common star-schema relationship is:

\begin{lstlisting}
Dimension (1) -> (*) Fact
\end{lstlisting}

For example:

\begin{lstlisting}
DimCustomer[CustomerID]
        |
        | 1:*
        |
FactSales[CustomerID]
\end{lstlisting}

Understand:

\begin{itemize}[leftmargin=*]
    \item one-to-one;
    \item one-to-many;
    \item many-to-many;
    \item active versus inactive relationships;
    \item filter direction.
\end{itemize}

Many-to-many and bidirectional relationships can be useful, but they should not
be the default solution to a poorly designed model.

\textbf{Practical rule}

Prefer simple:

\begin{lstlisting}
Dimension -> Fact
\end{lstlisting}

filtering whenever the model permits it.


% ------------------------------------------------------------
% DAX
% ------------------------------------------------------------

\textbf{DAX}

DAX (Data Analysis Expressions) is Power BI's analytical formula language.

A basic measure:

\begin{lstlisting}
Total Sales =
SUM(Sales[Revenue])
\end{lstlisting}

Average:

\begin{lstlisting}
Average Sale =
AVERAGE(Sales[Revenue])
\end{lstlisting}

Count rows:

\begin{lstlisting}
Number of Sales =
COUNTROWS(Sales)
\end{lstlisting}

Distinct count:

\begin{lstlisting}
Customers =
DISTINCTCOUNT(Sales[CustomerID])
\end{lstlisting}


% ------------------------------------------------------------
% MEASURES VS CALCULATED COLUMNS
% ------------------------------------------------------------

\textbf{Measures versus calculated columns}

This distinction is fundamental.

\begin{center}
\begin{tabularx}{\textwidth}{lXX}
\toprule
 & Measure & Calculated Column \\
\midrule
Computed & At query/report time & For each row \\
Depends on filters & Yes & Not in the same way \\
Stored per row & No & Yes \\
Typical use & Aggregation/KPI & Row-level derived value \\
\bottomrule
\end{tabularx}
\end{center}

Calculated column:

\begin{lstlisting}
Profit =
Sales[Revenue] - Sales[Cost]
\end{lstlisting}

Measure:

\begin{lstlisting}
Total Profit =
SUM(Sales[Revenue]) -
SUM(Sales[Cost])
\end{lstlisting}

For report-level aggregations, measures are generally preferable.


% ------------------------------------------------------------
% FILTER CONTEXT
% ------------------------------------------------------------

\textbf{Filter context}

Filter context is one of the most important concepts in DAX.

Suppose:

\begin{lstlisting}
Total Sales =
SUM(Sales[Revenue])
\end{lstlisting}

The same measure may display different values when the report is filtered by:

\begin{itemize}[leftmargin=*]
    \item year;
    \item product;
    \item region;
    \item customer;
    \item slicer selections.
\end{itemize}

The formula does not change. The \textbf{filter context} changes.

This is why measures dynamically respond to report interactions.


% ------------------------------------------------------------
% CALCULATE
% ------------------------------------------------------------

\textbf{CALCULATE}

\texttt{CALCULATE} evaluates an expression under modified filter context.

Example:

\begin{lstlisting}
US Sales =
CALCULATE(
    [Total Sales],
    Customer[Country] = "USA"
)
\end{lstlisting}

Another example:

\begin{lstlisting}
Sales 2026 =
CALCULATE(
    [Total Sales],
    DimDate[Year] = 2026
)
\end{lstlisting}

A useful way to remember it is:

\begin{quote}
\texttt{CALCULATE} changes the filters under which an expression is evaluated.
\end{quote}

This is one of the most important DAX concepts for interviews.


% ------------------------------------------------------------
% DIVIDE
% ------------------------------------------------------------

\textbf{DIVIDE}

Instead of:

\begin{lstlisting}
Profit Margin =
[Total Profit] / [Total Sales]
\end{lstlisting}

prefer:

\begin{lstlisting}
Profit Margin =
DIVIDE(
    [Total Profit],
    [Total Sales]
)
\end{lstlisting}

\texttt{DIVIDE} handles zero denominators more safely.


% ------------------------------------------------------------
% FILTER REMOVAL
% ------------------------------------------------------------

\textbf{Removing filters}

Example percentage of total:

\begin{lstlisting}
Percent of Total =
DIVIDE(
    [Total Sales],
    CALCULATE(
        [Total Sales],
        REMOVEFILTERS(Product)
    )
)
\end{lstlisting}

The denominator removes the current product filter while retaining other
relevant context.


% ------------------------------------------------------------
% ITERATOR FUNCTIONS
% ------------------------------------------------------------

\textbf{Iterator functions}

Functions ending in \texttt{X} commonly iterate over a table.

Examples:

\begin{lstlisting}
SUMX
AVERAGEX
MINX
MAXX
\end{lstlisting}

Example:

\begin{lstlisting}
Total Profit =
SUMX(
    Sales,
    Sales[Quantity] *
    (Sales[Price] - Sales[UnitCost])
)
\end{lstlisting}

\texttt{SUM} aggregates an existing column.

\texttt{SUMX} evaluates an expression row-by-row and then sums the results.


% ------------------------------------------------------------
% DATE TABLES
% ------------------------------------------------------------

\textbf{Date tables}

Time-based reporting should generally use a dedicated date dimension.

Useful fields include:

\begin{lstlisting}
Date
Year
Quarter
Month
Month Number
Week
Day
\end{lstlisting}

Connect the date table to the fact table using the appropriate date key.

A dedicated date table makes time-based filtering and comparisons much easier
to manage consistently.


% ------------------------------------------------------------
% VISUALS
% ------------------------------------------------------------

\textbf{Common visuals}

\begin{center}
\begin{tabularx}{\textwidth}{lX}
\toprule
Visual & Good Use \\
\midrule
Card & Single KPI \\
Bar Chart & Compare categories \\
Line Chart & Time series/trends \\
Scatter Plot & Relationship between numeric variables \\
Table & Detailed values \\
Matrix & Hierarchical/tabular summaries \\
Map & Geographic patterns \\
Slicer & Interactive filtering \\
\bottomrule
\end{tabularx}
\end{center}

Choose the visual based on the analytical question rather than visual novelty.


% ------------------------------------------------------------
% INTERACTIONS
% ------------------------------------------------------------

\textbf{Report interactions}

Power BI visuals can interact with one another through filtering and
highlighting.

A user selecting a category in one chart may change several other visuals on
the page.

Always test these interactions. A report can display individually correct
charts while producing confusing behavior when selections propagate across the
page.


% ------------------------------------------------------------
% IMPORT VS DIRECTQUERY
% ------------------------------------------------------------

\textbf{Import versus DirectQuery}

\begin{center}
\begin{tabularx}{\textwidth}{lXX}
\toprule
 & Import & DirectQuery \\
\midrule
Data location & Loaded into Power BI & Remains in source \\
Typical performance & Faster & Source dependent \\
Refresh needed & Yes & Queries source \\
Model flexibility & High & More constrained \\
Source load & Lower during viewing & Queries source during use \\
\bottomrule
\end{tabularx}
\end{center}

Import mode is often preferable when the dataset fits comfortably and
scheduled refreshes are acceptable.

DirectQuery is useful when querying the underlying source is necessary, but
report performance becomes more dependent on source-system performance and
model design.


% ------------------------------------------------------------
% PERFORMANCE
% ------------------------------------------------------------

\textbf{Performance tips}

\begin{itemize}[leftmargin=*]
    \item Prefer a clean star schema.
    \item Remove unused columns before loading large datasets.
    \item Avoid unnecessary high-cardinality text columns.
    \item Prefer measures for report-level calculations.
    \item Avoid excessive bidirectional relationships.
    \item Avoid unnecessarily complicated DAX when the model can be simplified.
    \item Reduce the number of unnecessary visuals on a page.
    \item Push appropriate transformations upstream when practical.
\end{itemize}


% ------------------------------------------------------------
% PRACTICAL TIPS
% ------------------------------------------------------------

\textbf{Practical tips}

\begin{itemize}[leftmargin=*]
    \item Determine table grain before creating relationships.
    \item Build the model before spending significant time formatting visuals.
    \item Prefer star schemas over one giant flattened table when appropriate.
    \item Use Power Query for data preparation and DAX for analytical logic.
    \item Prefer measures over calculated columns for dynamic aggregations.
    \item Learn \texttt{CALCULATE}; it is central to serious DAX work.
    \item Create a proper date table for time-based reporting.
    \item Test slicers and cross-filtering interactions.
    \item Use clear business names for measures.
    \item Validate important dashboard totals against the source data.
    \item Avoid putting every available metric on one page.
    \item Design each report page around a small number of analytical questions.
\end{itemize}


% ------------------------------------------------------------
% INTERVIEW NOTES
% ------------------------------------------------------------

\textbf{Interview notes}

Be comfortable explaining:

\begin{itemize}[leftmargin=*]
    \item Power Query versus DAX;
    \item fact tables versus dimension tables;
    \item star schemas;
    \item one-to-many relationships;
    \item measures versus calculated columns;
    \item row context versus filter context;
    \item the purpose of \texttt{CALCULATE};
    \item Import versus DirectQuery;
    \item why model design affects report performance.
\end{itemize}

A common interview question is:

\begin{quote}
Why is my Power BI measure returning different values in different visuals?
\end{quote}

A strong first answer is to investigate \textbf{filter context}: slicers,
visual-level filters, page filters, relationships, and the dimensions used by
each visual can change the context in which the measure is evaluated.


% ============================================================
% TABLEAU
% ============================================================

\subsection{Tableau}

\textbf{Tableau} is a visual analytics and business intelligence platform used
to explore data and build interactive dashboards.

The basic workflow is:

\begin{lstlisting}
Connect to Data
-> Define Fields
-> Build Worksheets
-> Add Filters / Calculations
-> Combine into Dashboard
\end{lstlisting}

Tableau is especially centered around interactive visual exploration and the
relationship between \textbf{dimensions}, \textbf{measures}, and
\textbf{level of detail}.


% ------------------------------------------------------------
% CORE TERMINOLOGY
% ------------------------------------------------------------

\textbf{Core terminology}

\begin{center}
\begin{tabularx}{\textwidth}{lX}
\toprule
Term & Meaning \\
\midrule
Worksheet & Individual visualization \\
Dashboard & Collection of worksheets and interactive objects \\
Story & Sequence of sheets/dashboards used to communicate a narrative \\
Dimension & Field generally used to categorize or segment data \\
Measure & Field generally used for quantitative aggregation \\
Discrete & Produces distinct headers/categories \\
Continuous & Produces a continuous axis \\
Calculated Field & Derived value defined using a Tableau expression \\
Parameter & User-controlled value usable in calculations and interactions \\
\bottomrule
\end{tabularx}
\end{center}


% ------------------------------------------------------------
% DIMENSIONS AND MEASURES
% ------------------------------------------------------------

\textbf{Dimensions versus measures}

Examples of dimensions:

\begin{lstlisting}
Customer
Region
Product
Category
Date
\end{lstlisting}

Examples of measures:

\begin{lstlisting}
Sales
Profit
Quantity
Cost
\end{lstlisting}

A simple way to think about them is:

\begin{itemize}[leftmargin=*]
    \item \textbf{Dimension}: how should the data be grouped?
    \item \textbf{Measure}: what should be measured or aggregated?
\end{itemize}

For example:

\begin{lstlisting}
Sales by Region
\end{lstlisting}

uses:

\begin{lstlisting}
Region -> Dimension
Sales  -> Measure
\end{lstlisting}


% ------------------------------------------------------------
% DISCRETE AND CONTINUOUS
% ------------------------------------------------------------

\textbf{Discrete versus continuous}

Tableau also distinguishes between discrete and continuous fields.

\begin{itemize}[leftmargin=*]
    \item \textbf{Discrete}: creates headers or categories.
    \item \textbf{Continuous}: creates an axis over a continuous range.
\end{itemize}

This distinction is separate from dimension versus measure.

A date, for example, can be displayed as either discrete categories or a
continuous timeline depending on the desired visualization.


% ------------------------------------------------------------
% ROWS AND COLUMNS
% ------------------------------------------------------------

\textbf{Rows and Columns shelves}

A simple bar chart might use:

\begin{lstlisting}
Columns: Region
Rows:    SUM(Sales)
\end{lstlisting}

A time series might use:

\begin{lstlisting}
Columns: Date
Rows:    SUM(Sales)
\end{lstlisting}

Dragging fields between Rows and Columns changes the structure of the
visualization.


% ------------------------------------------------------------
% MARKS CARD
% ------------------------------------------------------------

\textbf{Marks card}

The Marks card controls how observations are represented.

Common properties include:

\begin{itemize}[leftmargin=*]
    \item Color;
    \item Size;
    \item Label;
    \item Detail;
    \item Tooltip;
    \item Shape.
\end{itemize}

For example:

\begin{lstlisting}
Columns: Sales
Rows:    Profit
Color:   Region
Detail:  Customer
\end{lstlisting}

creates a scatter plot where customer observations are colored by region.


% ------------------------------------------------------------
% AGGREGATION
% ------------------------------------------------------------

\textbf{Aggregation}

Measures are usually aggregated.

Common aggregations include:

\begin{lstlisting}
SUM(Sales)
AVG(Sales)
MIN(Sales)
MAX(Sales)
COUNT(Sales)
COUNTD(CustomerID)
\end{lstlisting}

\texttt{COUNTD} means distinct count.

Always check the aggregation being used. A visualization of:

\begin{lstlisting}
SUM(Profit)
\end{lstlisting}

answers a different question from:

\begin{lstlisting}
AVG(Profit)
\end{lstlisting}


% ------------------------------------------------------------
% CALCULATED FIELDS
% ------------------------------------------------------------

\textbf{Calculated fields}

Basic arithmetic:

\begin{lstlisting}
[Sales] - [Cost]
\end{lstlisting}

Conditional logic:

\begin{lstlisting}
IF [Profit] > 0 THEN
    "Profitable"
ELSE
    "Unprofitable"
END
\end{lstlisting}

Multiple conditions:

\begin{lstlisting}
IF [Sales] >= 10000 THEN
    "High"
ELSEIF [Sales] >= 5000 THEN
    "Medium"
ELSE
    "Low"
END
\end{lstlisting}

Handle missing values:

\begin{lstlisting}
IFNULL([Profit], 0)
\end{lstlisting}


% ------------------------------------------------------------
% TABLE CALCULATIONS
% ------------------------------------------------------------

\textbf{Table calculations}

Table calculations operate over values already present in the visualization.

Examples include:

\begin{itemize}[leftmargin=*]
    \item running totals;
    \item percent of total;
    \item rank;
    \item moving averages;
    \item differences from previous values.
\end{itemize}

Running total:

\begin{lstlisting}
RUNNING_SUM(SUM([Sales]))
\end{lstlisting}

Rank:

\begin{lstlisting}
RANK(SUM([Sales]))
\end{lstlisting}

Previous value:

\begin{lstlisting}
LOOKUP(SUM([Sales]), -1)
\end{lstlisting}

Table calculations depend heavily on how the visualization is partitioned and
addressed. If a table calculation produces unexpected results, inspect the
direction and level over which it is computing.


% ------------------------------------------------------------
% LEVEL OF DETAIL EXPRESSIONS
% ------------------------------------------------------------

\textbf{Level of Detail (LOD) expressions}

LOD expressions explicitly control the granularity at which a calculation is
performed.

\textbf{FIXED}

Calculate customer-level sales regardless of the dimensions currently shown in
the view:

\begin{lstlisting}
{ FIXED [Customer ID] :
    SUM([Sales])
}
\end{lstlisting}

\textbf{INCLUDE}

Include an additional dimension in the calculation:

\begin{lstlisting}
{ INCLUDE [Customer ID] :
    SUM([Sales])
}
\end{lstlisting}

\textbf{EXCLUDE}

Remove a dimension from the calculation:

\begin{lstlisting}
{ EXCLUDE [Region] :
    SUM([Sales])
}
\end{lstlisting}

The three major LOD expressions are:

\begin{itemize}[leftmargin=*]
    \item \texttt{FIXED}: calculate at an explicitly specified level.
    \item \texttt{INCLUDE}: add dimensions to the calculation's level.
    \item \texttt{EXCLUDE}: remove dimensions from the calculation's level.
\end{itemize}

LOD expressions are one of the most important Tableau concepts for technical
interviews.


% ------------------------------------------------------------
% FILTERS
% ------------------------------------------------------------

\textbf{Filters}

Common filters include:

\begin{itemize}[leftmargin=*]
    \item data source filters;
    \item extract filters;
    \item context filters;
    \item dimension filters;
    \item measure filters;
    \item table-calculation filters.
\end{itemize}

Filtering order can matter, especially when LOD expressions or Top-N filters
are involved.

If a calculation produces an unexpected result, consider whether the relevant
filter is applied before or after the calculation.


% ------------------------------------------------------------
% CONTEXT FILTERS
% ------------------------------------------------------------

\textbf{Context filters}

A context filter creates a subset of data that other filters can operate on.

This is particularly useful when a secondary filter should be evaluated only
within an already filtered subset.

For example:

\begin{lstlisting}
Context Filter: Region = East
Top-N Filter:   Top 10 Customers
\end{lstlisting}

can mean:

\begin{lstlisting}
Top 10 customers within the East region
\end{lstlisting}

rather than top customers globally followed by an East-region filter.


% ------------------------------------------------------------
% PARAMETERS
% ------------------------------------------------------------

\textbf{Parameters}

Parameters are user-controlled values that can drive calculations.

For example, create a parameter:

\begin{lstlisting}
Top N
\end{lstlisting}

and use it in a calculated field:

\begin{lstlisting}
RANK(SUM([Sales])) <= [Top N]
\end{lstlisting}

Parameters can support:

\begin{itemize}[leftmargin=*]
    \item Top-N controls;
    \item thresholds;
    \item scenario analysis;
    \item metric selection;
    \item reference values.
\end{itemize}


% ------------------------------------------------------------
% RELATIONSHIPS AND JOINS
% ------------------------------------------------------------

\textbf{Relationships versus joins}

Tableau can combine tables using relationships or physical joins.

A physical join combines rows before visualization:

\begin{lstlisting}
Orders
LEFT JOIN
Customers
\end{lstlisting}

Relationships allow tables to remain logically distinct and can preserve their
individual levels of detail until they are needed in a visualization.

As with SQL, always understand:

\begin{lstlisting}
1:1
1:N
N:1
N:M
\end{lstlisting}

before combining tables.

Unexpected row multiplication can produce incorrect measures.


% ------------------------------------------------------------
% LIVE CONNECTIONS AND EXTRACTS
% ------------------------------------------------------------

\textbf{Live connections versus extracts}

\begin{center}
\begin{tabularx}{\textwidth}{lXX}
\toprule
 & Live & Extract \\
\midrule
Data location & Source system & Tableau extract \\
Freshness & Current source data & Refresh dependent \\
Performance & Source dependent & Often faster \\
Offline use & Limited & Possible \\
Source load & Queries source & Reduced during analysis \\
\bottomrule
\end{tabularx}
\end{center}

An extract is a stored analytical snapshot optimized for Tableau.

Use the approach appropriate for the required balance between freshness and
performance.


% ------------------------------------------------------------
% DASHBOARDS
% ------------------------------------------------------------

\textbf{Dashboards}

Dashboards combine worksheets into an interactive reporting interface.

Useful dashboard components include:

\begin{itemize}[leftmargin=*]
    \item worksheets;
    \item filters;
    \item legends;
    \item parameters;
    \item text;
    \item navigation controls.
\end{itemize}

A useful dashboard should answer a clear analytical question rather than simply
display as many charts as possible.


% ------------------------------------------------------------
% DASHBOARD ACTIONS
% ------------------------------------------------------------

\textbf{Dashboard actions}

Actions allow one visualization to affect another.

Common actions include:

\begin{itemize}[leftmargin=*]
    \item filter actions;
    \item highlight actions;
    \item parameter actions;
    \item navigation actions.
\end{itemize}

For example, selecting a state on a map can filter a sales chart to that state.

Test interactions carefully so users understand what is currently filtering
the dashboard.


% ------------------------------------------------------------
% VISUAL SELECTION
% ------------------------------------------------------------

\textbf{Common visual choices}

\begin{center}
\begin{tabularx}{\textwidth}{lX}
\toprule
Visual & Good Use \\
\midrule
Bar Chart & Category comparisons \\
Line Chart & Trends over time \\
Scatter Plot & Relationships between numeric variables \\
Map & Geographic patterns \\
Histogram & Distribution of a numeric variable \\
Box Plot & Distribution and group comparison \\
Heat Map & Patterns across two categorical dimensions \\
Table & Detailed values \\
\bottomrule
\end{tabularx}
\end{center}

Avoid choosing a visualization simply because it looks interesting. The
visual should make the intended comparison easier to understand.


% ------------------------------------------------------------
% PERFORMANCE
% ------------------------------------------------------------

\textbf{Performance tips}

\begin{itemize}[leftmargin=*]
    \item Avoid loading unnecessary columns and rows.
    \item Reduce extremely high-cardinality dimensions when they are not needed.
    \item Avoid excessive marks in a single visualization.
    \item Use extracts when appropriate for performance.
    \item Avoid unnecessary complex calculations.
    \item Understand the grain of joined or related tables.
    \item Reduce unnecessary dashboard objects.
\end{itemize}


% ------------------------------------------------------------
% PRACTICAL TIPS
% ------------------------------------------------------------

\textbf{Practical tips}

\begin{itemize}[leftmargin=*]
    \item Always check whether a field is a dimension or measure.
    \item Check whether a field is discrete or continuous.
    \item Verify the aggregation applied to every important measure.
    \item Understand the data grain before combining tables.
    \item Use LOD expressions when a calculation needs a specific granularity.
    \item Check table-calculation direction when results appear incorrect.
    \item Use parameters for user-controlled analytical behavior.
    \item Keep dashboards focused on a small number of questions.
    \item Test filter and dashboard actions before publishing.
    \item Validate important totals against the source data.
    \item Use informative titles that describe what a visual actually shows.
\end{itemize}


% ------------------------------------------------------------
% INTERVIEW NOTES
% ------------------------------------------------------------

\textbf{Interview notes}

Be comfortable explaining:

\begin{itemize}[leftmargin=*]
    \item dimensions versus measures;
    \item discrete versus continuous fields;
    \item aggregation;
    \item calculated fields;
    \item table calculations;
    \item LOD expressions;
    \item \texttt{FIXED}, \texttt{INCLUDE}, and \texttt{EXCLUDE};
    \item context filters;
    \item parameters;
    \item relationships versus joins;
    \item live connections versus extracts.
\end{itemize}

A common interview question is the difference between a table calculation and
an LOD expression.

A useful answer is:

\begin{itemize}[leftmargin=*]
    \item A \textbf{table calculation} operates on the aggregated values present in
    the visualization and depends on the layout of the view.
    \item An \textbf{LOD expression} explicitly controls the granularity at which a
    calculation is performed relative to the visualization.
\end{itemize}

Another common question is when to use \texttt{FIXED}. Use it when a metric
must be calculated at a specific level of detail regardless of some dimensions
currently present in the view.


% ============================================================
% POWER BI VS TABLEAU
% ============================================================

\textbf{Power BI versus Tableau}

\begin{center}
\begin{tabularx}{\textwidth}{lXX}
\toprule
Concept & Power BI & Tableau \\
\midrule
Data transformation & Power Query & Data Source / Prep workflow \\
Calculation language & DAX & Tableau calculations \\
Core model concept & Semantic/star model & Dimensions/measures and LOD \\
Reusable aggregation & Measure & Calculated field / LOD \\
Granularity control & Filter context / DAX & LOD expressions \\
User input & Parameters / slicers & Parameters / filters \\
Cached data & Import model & Extract \\
Source querying & DirectQuery & Live connection \\
Primary report unit & Report page & Worksheet/dashboard \\
\bottomrule
\end{tabularx}
\end{center}

For interviews, avoid framing the difference as simply ``Power BI is Microsoft
and Tableau is Salesforce.'' The more useful technical distinction is how each
tool approaches modeling and calculations.

In Power BI, spend particular attention on:

\begin{lstlisting}
Data Model
Relationships
Measures
Filter Context
DAX
\end{lstlisting}

In Tableau, spend particular attention on:

\begin{lstlisting}
Dimensions / Measures
Discrete / Continuous
Level of Detail
Table Calculations
Filters
\end{lstlisting}

For both tools, the most important reporting principle is the same:

\begin{quote}
A visually polished dashboard is still wrong if the underlying data model,
aggregation, filtering, or grain is wrong.
\end{quote}


% ============================================================
% END DOCUMENT
% ============================================================

\end{document}